% concatenated from On_the_representability_of_actions_of_unital_algebras.tex
\documentclass[reqno,a4paper]{amsart}
\usepackage[T1]{fontenc}
\usepackage[utf8]{inputenc}
\usepackage[english]{babel}
\usepackage{microtype}

\newcommand{\lc}{$\ell$}
\usepackage{latexsym, amssymb, amsmath}
\usepackage{enumitem}
\usepackage{amsthm}
\usepackage{tikz-cd}
\usepackage{orcidlink}
\usepackage{doi}
\renewcommand{\doi}[1]{DOI:~\href{https://doi.org/#1}{#1}}
\usepackage{url}

%\usepackage[british]{babel}
\usepackage{verbatim}
\usepackage[all]{xy}
\usepackage{calrsfs}
\usepackage{mathabx}
\usepackage{xcolor}
\usepackage{mathtools}
\usepackage{blkarray}
\usepackage[textsize=scriptsize]{todonotes}
\usepackage{thmtools}
\usepackage{xurl}
\usepackage{nameref,hyperref,cleveref}
\usepackage{quiver}

\theoremstyle{plain}
\newtheorem{theorem}[subsection]{Theorem}
\newtheorem{proposition}[subsection]{Proposition}
\newtheorem{corollary}[subsection]{Corollary}	
\newtheorem{lemma}[subsection]{Lemma}

\theoremstyle{definition}
\newtheorem{definition}[subsection]{Definition}

\theoremstyle{remark}
\newtheorem{remark}[subsection]{Remark}
\newtheorem{example}[subsection]{Example}
\newtheorem{examples}[subsection]{Examples}
\newtheorem{problem}[subsection]{Open problem}

\setcounter{tocdepth}{1}
\numberwithin{equation}{section}

\newenvironment{tfae}
{\renewcommand{\labelenumi}{(\roman{enumi})}
\begin{enumerate}}
{\end{enumerate}}
\renewcommand{\labelitemi}{--}

\newcommand{\defn}{\textbf}
\newcommand{\noproof}{\hfil\qed}
\newcommand{\cd}{\cdot}
\newcommand{\ra}{\rightarrow}
\newcommand{\ot}{\leftarrow}
\newcommand{\tensor}{\otimes}

\newcommand{\C}{\ensuremath{\mathcal{C}}}
\newcommand{\D}{\ensuremath{\mathcal{D}}}
\newcommand{\A}{\ensuremath{\mathsf{A}}}
\newcommand{\B}{\ensuremath{\mathsf{B}}}
\newcommand{\Lr}{\ensuremath{\mathcal{L}}}
\newcommand{\Mr}{\ensuremath{\mathcal{M}}}
\newcommand{\Part}{\ensuremath{\mathcal{P}}}
\newcommand{\F}{\mathbb{F}}         
%\newcommand{\cV}{\mathsf{V}}
\newcommand{\E}{\mathcal{E}}
\newcommand{\K}{\mathbb{K}}
\newcommand{\N}{\mathbb{N}}
\newcommand{\V}{\mathcal{V}}
\newcommand{\Z}{\mathbb{Z}}
\newcommand{\U}{\mathcal{U}}
\newcommand{\mh}{\mathfrak{h}}
\newcommand{\mg}{\mathfrak{g}}

\DeclareMathOperator{\bchar}{char}
\DeclareMathOperator{\Ker}{Ker}
\DeclareMathOperator{\Aut}{Aut}
\DeclareMathOperator{\op}{op}
\DeclareMathOperator{\Act}{Act}
\DeclareMathOperator{\Hom}{Hom}
\DeclareMathOperator{\Bim}{Bim}
\DeclareMathOperator{\SplExt}{SplExt}
\DeclareMathOperator{\Imm}{Im}
\DeclareMathOperator{\End}{End}
\DeclareMathOperator{\Der}{Der}
\DeclareMathOperator{\ADer}{ADer}
\DeclareMathOperator{\Bider}{Bider}
\DeclareMathOperator{\Inn}{Inn}
\DeclareMathOperator{\M}{M}
\DeclareMathOperator{\bZ}{Z}
\DeclareMathOperator{\USGA}{USGA}
\DeclareMathOperator{\Ann}{Ann}
\DeclareMathOperator{\FM}{FM}
\DeclareMathOperator{\Ob}{Ob}
\DeclareMathOperator{\Gal}{Gal}
\DeclareMathOperator{\Fam}{Fam}
\DeclareMathOperator{\id}{id}

\newcommand{\Hp}{\ensuremath{\mathsf{Hoops}}}
\newcommand{\WH}{\ensuremath{\mathsf{WHoops}}}
\newcommand{\PH}{\ensuremath{\mathsf{PHoops}}}
\newcommand{\MV}{\ensuremath{\mathsf{MVAlg}}}
\newcommand{\PA}{\ensuremath{\mathsf{PrAlg}}}
\newcommand{\Gp}{\ensuremath{\mathsf{Grp}}}
\newcommand{\GSet}{\ensuremath{G\text{-}\mathsf{Set}}}
\newcommand{\Set}{\ensuremath{\mathsf{Set}}}
\newcommand{\Rng}{\ensuremath{\mathsf{Rng}}}
\newcommand{\Mag}{\ensuremath{\mathsf{Mag}}}
\newcommand{\Pt}{\ensuremath{\mathsf{Pt}}}
\newcommand{\Assoc}{\ensuremath{\mathsf{Assoc}}}
\newcommand{\AAssoc}{\ensuremath{\mathsf{AAssoc}}}
\newcommand{\Lie}{\ensuremath{\mathsf{Lie}}}
\newcommand{\Leib}{\ensuremath{\mathsf{Leib}}}
\newcommand{\Pois}{\ensuremath{\mathsf{Pois}}}
\newcommand{\CPois}{\ensuremath{\mathsf{CPois}}}
\newcommand{\BoolRng}{\ensuremath{\mathsf{BoolRng}}}
\newcommand{\BoolAlg}{\ensuremath{\mathsf{BoolAlg}}}
\newcommand{\CAssoc}{\ensuremath{\mathsf{CAssoc}}}
\newcommand{\ACAAssoc}{\ensuremath{\mathsf{ACAAssoc}}}
\newcommand{\Com}{\ensuremath{\mathsf{Com}}}
\newcommand{\ACom}{\ensuremath{\mathsf{ACom}}}
\newcommand{\JJord}{\ensuremath{\mathsf{JJord}}}
\newcommand{\Jord}{\ensuremath{\mathsf{Jord}}}
\newcommand{\AbAlg}{\ensuremath{\mathsf{AbAlg}}}
\newcommand{\Alg}{\ensuremath{\mathsf{Alg}}}
\newcommand{\PAlg}{\ensuremath{\mathsf{PAlg}}}
\newcommand{\Nil}{\ensuremath{\mathsf{Nil}}}
\newcommand{\Sol}{\ensuremath{\mathsf{Sol}}}
\newcommand{\Alt}{\ensuremath{\mathsf{Alt}}}
\newcommand{\WAlt}{\ensuremath{\mathsf{WAlt}}}
\newcommand{\Grp}{\ensuremath{\mathsf{Grp}}}
\newcommand{\Vect}{\ensuremath{\mathsf{Vec}}}
\newcommand{\Cat}{\ensuremath{\mathsf{Cat}}}
\newcommand{\Ring}{\ensuremath{\mathsf{Ring}}}
\newcommand{\CRng}{\ensuremath{\mathsf{CRng}}}
\newcommand{\CRing}{\ensuremath{\mathsf{CRing}}}
\newcommand{\SSS}{\ensuremath{\mathbb{S}}}
\newcommand{\SplExtX}{\ensuremath{\mathsf{SplExt}(X)}}

\newdir{>}{{}*:(1,-.2)@^{>}*:(1,+.2)@_{>}}
\newdir{<}{{}*:(1,+.2)@^{<}*:(1,-.2)@_{<}}
\newdir{>>}{{}*!/3.5pt/:(1,-.2)@^{>}*!/3.5pt/:(1,+.2)@_{>}*!/7pt/:(1,-.2)@^{>}*!/7pt/:(1,+.2)@_{>}}
\newdir{ >}{{}*!/-8pt/@{>}}

\date{}
\newcommand{\ulm}{\mathsf{u}\text{\lc}\mathsf{M}}
\newcommand{\ulmp}{\mathsf{u}\text{\lc}\mathsf{M}^{+}}
\newcommand{\ua}{\underline{\alpha}}
\newcommand{\ux}{\underline{x}}
\newcommand{\w}{\rightarrow}
\newcommand{\pt}{\operatorname{Pt}}

\newsavebox{\pullback}
\sbox\pullback{%
\begin{tikzpicture}%
\draw (0,0) -- (1ex,0ex);%
\draw (1ex,0ex) -- (1ex,1ex);%
\end{tikzpicture}}

\begin{document}

\title[On the representability of actions of unital algebras]{On the representability of actions\\of unital algebras}

%\author[M.~Mancini]{M.~Mancini$^{*}$~\orcidlink{0000-0003-2142-6193}}
\author[M.~Mancini]{Manuel Mancini~\orcidlink{0000-0003-2142-6193}}
\author[F.~Piazza]{Federica Piazza~\orcidlink{0009-0001-1028-9659}}
\author[C.~Vienne]{Corentin Vienne~\orcidlink{0000-0002-0397-3495}}

\email{manuel.mancini@uclouvain.be; manuel.mancini@unipa.it}
\email{federica.piazza07@unipa.it; federica.piazza1@studenti.unime.it}
\email{corentin.vienne@proton.me}

\address[M.~Mancini]{Institut de Recherche en Mathématique et Physique, Université catholique de Louvain, chemin du cyclotron 2 bte L7.01.02, B--1348 Louvain-la-Neuve, Belgium.}
\address[M.~Mancini, F.~Piazza]{Dipartimento di Matematica e Informatica, Università degli Studi di Palermo, Via Archirafi 34, 90123 Palermo, Italy.}
\address[F.~Piazza]{Dipartimento di Scienze Matematiche e Informatiche, Scienze Fisiche e Scienze della Terra, Università degli Studi di Messina, Viale Ferdinando Stagno d'Alcontres 31, 98166 Messina, Italy.}
\address[C.~Vienne]{Dipartimento di Scienze per gli Alimenti, la Nutrizione, e l'Ambiente, Università degli Studi di Milano Statale, Via Celoria 2, 20133 Milano, Italy.}

%\thanks{$^*$Corresponding Author.}

\begin{abstract}
Working in the setting of \emph{ideally exact categories}, we investigate the representability of actions of unital non-associative algebras over a field. We show that, in general, such categories fail to be \emph{action representable}: for instance, the category of all unital algebras is not even \emph{action accessible}.

We then consider this problem in the context of operadic, action accessible, unit-closed varieties. Using the construction of the \emph{external weak actor}, we prove that for any algebra $X$ in such a variety $\V$, the canonical map into its external weak actor is an isomorphism if and only if $X$ is unital. Consequently, the ideally exact category $\V_1$ of unital algebras in $\V$ is action representable, and the actor of $X$ is $X$ itself.

Finally, we prove action representability for unital Poisson algebras via an explicit construction of the \emph{universal strict general actor}.
\end{abstract}

\subjclass[2020]{08A35; 08C05; 16B50; 16W25; 17A36; 17B63; 18E13}
\keywords{Action representable category, ideally exact category, internal action, split extension, non-associative algebra, unital algebra, Poisson algebra}

\maketitle

\section{Introduction}\label{SecIntro}

The concept of \emph{internal object action} was introduced by F.~Borceux, G.~Janelidze, and G.~M.~Kelly in~\cite{BJK} with the goal of extending the correspondence between actions and split extensions from the setting of groups to the more general setting of semi-abelian categories~\cite{Semi-Ab}. In some of those categories, internal actions are exceptionally well behaved, in the sense that the actions on each object~$X$ are \emph{representable}: this means that there exists an object~$[X]$, called the \emph{actor} of $X$, such that the functor $\Act(-,X) \cong \SplExt(-,X)$ which sends an object $B$ to the set of internal actions/isomorphism classes of split extensions of $B$ on/by $X$, is naturally isomorphic to the functor $\Hom(-,[X])$. This is the case for the category of groups, where $[X]=\Aut(X)$ is the group of automorphisms on $X$. 

The study of action representability in semi-abelian categories was further developed in~\cite{BJK2}, where it was shown, for example, that the variety of commutative associative algebras over a field fails to be action representable. Later, the article~\cite{Tim} established that amongst varieties of non-associative algebras over an infinite field of characteristic different from $2$, only the category of abelian algebras and that of Lie algebras satisfy this property. The restrictive nature of action representability naturally led to the introduction of weaker, yet related, notions.

In~\cite{act_accessible}, D.~Bourn and G.~Janelidze introduced the notion of \emph{action accessibility} to encompass significant examples of categories that do not satisfy action representability, such as (not necessarily unital) rings, associative algebras, and Leibniz algebras~\cite{loday}. A.~Montoli later proved in~\cite{Montoli} that every \emph{Orzech category of interest}~\cite{Orzech} is action accessible. Additionally, the authors of~\cite{Casas} presented a broader notion of representing object in any Orzech category of interest: the \emph{universal strict general actor}.

More recently, G.~Janelidze introduced in~\cite{WRA} the notion of \emph{weak action representability}. This condition requires for each object $X$ of a semi-abelian category the existence of a \emph{weak representation} of the actions on $X$, i.e., an object $T(X)$ along with a natural monomorphism of functors $\SplExt(-,X) \rightarrowtail \Hom(-,T(X))$. Examples of weakly action representable categories are the variety of associative algebras~\cite{WRA} and the variety of Leibniz algebras~\cite{CigoliManciniMetere}.

In~\cite{WRAAlg}, the authors explored the concept of weak action representability within the framework of varieties of non-associative algebras over a field. They worked towards the construction of an external weakly representing object $\E(X)$ for actions on an object $X$ of a variety of non-associative algebras $\V$. They actually obtained a \emph{partial algebra} $\E(X)$, called \emph{external weak actor} of $X$, together with a monomorphism of functors $\SplExt(-,X) \rightarrowtail \Hom_\PAlg(\widetilde{U}(-),\E(X))$, where $\PAlg$ is the category of partial algebras and $\widetilde{U} \colon \V \to \PAlg$ denotes the forgetful functor.

G.~Janelidze later extended the notions of action accessibility and (weak) action representability to the broader setting of \emph{ideally exact categories}~\cite{IdeallyExact}, which were introduced as a generalization of semi-abelian categories in order to include relevant examples of non-pointed categories, such as the categories $\Ring$ and $\CRing$ of unital rings and unital commutative rings, all non-trivial protomodular varieties of universal algebras, and any cotopos. In this context, the action representability of an ideally exact category~$\U$ is defined relative to the semi-abelian category $(\U \downarrow 0)$, where~$0$ denotes the initial object of $\U$--see \Cref{def_act_repr}. It was also shown in~\cite{IdeallyExact} that both $\Ring$ and $\CRing$ are action representable in this broader sense.

The aim of this manuscript is to study the representability of actions within the setting of categories of unital non-associative algebras. After recalling the necessary background in Sections \ref{SecVar} and \ref{SecAct}, we show in \Cref{SecIdEx} that in general a category of unital algebras fails to be action representable. In fact, we prove that the categories of unital algebras and unital commutative algebras (both not necessarily associative) are not even action accessible. 

We then study in \Cref{SecMain} the representability of actions of unital algebras in any operadic, action accessible, unit-closed variety $\V$. Using the construction of the external weak actor $\E(X)$ of an algebra $X$ of $\V$, we prove that the canonical partial algebra homomorphism $\Inn \colon X \to \E(X)$, defined by $\Inn(x) = (L_x,R_x)$, where $L_x$ and $R_x$ denote the left and right multiplications by $x$, is an isomorphism if and only if~$X$ is unital. In \Cref{final_thm}, the main result of this article, we prove that the ideally exact category $\V_1$ of unital algebras in $\V$ is action representable, and the actor of a unital algebra $X$ is $X$ itself. This result implies that, in particular, the categories of unital alternative, unital associative, and unital commutative associative algebras are action representable, and highlights a contrast with the non-unital case, where action representability is a highly restrictive property, essentially characterising the variety of Lie algebras~\cite{Tim}.

Finally, in \Cref{SecPois} we study the representability of actions of the categories $\Pois_1$ and $\CPois_1$ of unital Poisson algebras and unital commutative Poisson algebras. Here we employ the explicit construction of the universal strict general actor given in~\cite{CigoliManciniMetere}. While it remains open whether the category of (all) Poisson algebras has weakly representable actions, we prove that the ideally exact categories $\Pois_1$ and $\CPois_1$ are action representable, and that the actor of a unital (commutative) Poisson algebra $X$ is isomorphic to $X$ itself.

\section{Varieties of non-associative algebras}\label{SecVar}

We now describe the algebraic framework in which we work: \emph{varieties of non-associative algebras} over a field $\F$. We think of those as collections of algebras satisfying a chosen set of polynomial equations. We refer the reader to~\cite{VdL-NAA} for more details.

A \emph{non-associative algebra} over $\F$ is a vector space $X$ equipped with a bilinear multiplication $X \times X \to  X$, denoted by $(x,y) \mapsto xy$. In general, the existence of a multiplicative identity is not assumed. The category of \emph{all} non-associative algebras over $\F$ is denoted by $\Alg$ and it has as morphisms the linear maps that preserve the multiplication.

\begin{definition}\label{def identity variety}
An \emph{identity} of an algebra $X$ is a non-associative polynomial
\[
\varphi=\varphi(x_1,\ldots,x_n)
\]
such that $\varphi(x_1,\ldots,x_n)=0$ for any $x_1,\ldots,x_n \in X$. We say that the algebra $X$ \emph{satisfies} the identity $\varphi$.
\end{definition}

If $I$ is a set of identities, then the \emph{variety of non-associative algebras} $\V$ determined by $I$ is the class of all algebras which satisfy all the identities in $I$. Conversely, we say that a variety of non-associative algebras $\V$ satisfies the identities in $I$ if every algebra of $\V$ satisfies the given set of identities. 

\begin{remark}
From now on, when we consider a variety $\V$, we assume that $\V$ is a variety of non-associative algebras over a field $\F$. We fix the field $\F$, so that we may drop it from our notation.
\end{remark}

For a monomial~$\varphi=\varphi(x_1,\ldots,x_n)$, we define its \emph{type} as the $n$-tuple $(k_1, \ldots , k_n)$, where $k_i$ is the number of times $x_i$ appears in $\varphi$, and its \emph{degree} as the natural number $k_1 + \cdots + k_n$. A polynomial is said to be \emph{homogeneous} if all its monomials are of the same type, and it is said to be \emph{multilinear} if all its monomials have type~$(1,\ldots,1)$.

\begin{definition}
A variety of non-associative algebras $\V$ is said to be \emph{homogeneous} (respectively \emph{operadic}) if it is determined by a set of homogeneous (respectively multilinear) identities.
\end{definition}

We recall that, if the field $\F$ is infinite, then any variety of non-associative $\F$-algebras $\V$ is homogeneous~\cite{hmg}. Furthermore, if $\mathrm{char}(\F)=0$, then any variety $\V$ is operadic. This follows from the multilinearisation process, see~\cite[Corollary 3.7]{Osborn}.

\begin{examples}\label{Examples varieties}{\ }
\begin{enumerate}
\item $\AbAlg$ is the variety of \emph{abelian} algebras, which is determined by the identity~$xy=0$.

\item $\Com$ and $\ACom$ denote, respectively, the varieties of commutative and anti-commutative (not necessarily associative) algebras. These varieties are determined, respectively, by the identities $xy-yx=0$ and $xy+yx=0$.

\item $\Assoc$ is the variety of \emph{associative} algebras, which is determined by the identity $x(yz)-(xy)z=0$.

\item $\CAssoc$ denotes the variety of commutative associative algebras.

\item $\Alt$ is the variety of \emph{alternative algebras}, which is determined by the identities
\[
(yx)x-y(xx)=0 \quad \text{and} \quad x(xy)-(xx)y=0.
\]
If $\bchar(\F) \neq 2$, then, as described in Example~\ref{example Alt}, the variety $\Alt$ can be interpreted as an operadic variety.

We observe that every associative algebra is alternative, while an example of an alternative algebra which is not associative is given by the algebra of \emph{octonions}.

\item $\Lie$ denotes the variety of \emph{Lie algebras}, characterised by $x^2=0$ together with the \emph{Jacobi identity} $x(yz)+y(zx)+z(xy)=0$. We recall that when $\bchar(\F) \neq 2$, the condition $x^2=0$ is equivalent to anti-commutativity, and hence in this situation the variety $\Lie$ is operadic.

\item $\JJord$ is the variety of \emph{Jacobi-Jordan algebras}~\cite{Burde}, which is determined by commutativity and the Jacobi identity. Jacobi–Jordan algebras, also known as \emph{mock-Lie algebras}, are the commutative counterpart of Lie algebras. 

\item $\Leib$ is the variety of \emph{(right) Leibniz algebras}, which is determined by the \emph{(right) Leibniz identity}, that is $(xy)z-(xz)y-x(yz)=0$.

\item $\BoolRng$ is the category of (not necessarily unital) \emph{Boolean rings}, which may be seen as the variety of commutative associative $\F_2$-algebras satisfying $x^2=x$. We remark that this variety is not homogeneous.

\item Given any variety $\V$, one can define the subvariety $\Nil_k(\V)$ of $k$-\emph{nilpotent} algebras of $\V$, which is determined by all the identities of $\V$ together with
\[
x_1 x_2 \cdots x_{k+1}=0,
\]
for all possible choices of parentheses. For instance, if $\V=\CAssoc$, every non-trivial multilinear identity of degree $k$ is equivalent to $x_1\cdots x_k=0$. Consequently, the only proper operadic subvarieties of $\CAssoc$ are the varieties of $k$-nilpotent commutative associative algebras.
\end{enumerate}
\end{examples}

In this paper, we aim to study the representability of actions of unital algebras. To do this, we need to recall the following definition.

\begin{definition}
A variety $\V$ is said to be \emph{unit-closed} if for any algebra $X$ of~$\V$, the algebra $\widetilde{X}=\langle X,1 \rangle$ obtained by adjoining to $X$ an external unit $1$ which satisfies the identities $x \cdot 1 = 1 \cdot x = x$ for any $x \in X$, is still an object of~$\V$. 
\end{definition}

For instance, the varieties $\Com$, $\Assoc$, $\CAssoc$, $\Alt$, and $\BoolRng$ are unit-closed, whereas the varieties $\AbAlg$, $\Leib$, and, more generally, any variety of anti-commutative algebras over a field of characteristic different from $2$, are not unit-closed. Note that in characteristic $2$, commutativity and anti-commutativity coincide. Thus, the property of being unit-closed depends on the identities satisfied by the variety, and the characteristic of the base field plays a crucial role. For example, the variety $\JJord$ is unit-closed if and only if $\operatorname{char}(\F)=3$. 

\begin{remark}
Given a unit-closed variety $\V$, we denote by $\V_1$ the subclass of unital algebras of $\V$. We will see in \Cref{SecIdEx} that $\V_1$ is an \emph{ideally exact category} in the sense of G.~Janelidze~\cite{IdeallyExact}.
\end{remark}

\section{Representability of actions in semi-abelian categories}\label{SecAct}

Semi-abelian categories were introduced in~\cite{Semi-Ab} by G.~Janelidze, L.~Márki and W.~Tholen in order to provide a categorical setting which would capture categorical-algebraic properties of groups, rings and algebras. We recall that a category~$\C$ is semi-abelian when it is pointed, Barr-exact~\cite{Barr}, Bourn-protomodular~\cite{Bourn} and has finite coproducts.

\begin{example}
Any variety of non-associative algebras $\V$ can be seen as a category where the morphisms are the same ones as in $\Alg$ and we have a full inclusion functor $\V \hookrightarrow \Alg$. In particular, any such variety is a semi-abelian category, see~\cite[Theorem 9.5]{VdL-NAA}.
\end{example}

A notion that can be explored in the context of semi-abelian categories is that of split extension.

\begin{definition}
A \emph{split extension} of an object $B$ by another object $X$ of a semi-abelian category $\C$ is a diagram
\begin{equation*}\label{eq:split_ext}
\begin{tikzcd}
X\arrow [r, "k"]
&A \arrow[r, shift left, "\alpha"] &
B\ar[l, shift left, "\beta"]
\end{tikzcd}
\end{equation*}
in $\C$, such that $\alpha \circ \beta = \id_B$ and $(X,k)$ is a kernel of $\alpha$.
\end{definition}

For any object $X$ of $\C$, one can consider the functor
\[
\SplExt(-,X) \colon \C^{\op} \to \Set
\]
which assigns to each object $B$ of $\C$ the set $\SplExt(B,X)$ of isomorphism classes of split extensions of $B$ by $X$, and to any morphism $f \colon B'\to  B$ the \emph{change of base} map $f^* \colon \SplExt(B,X) \to \SplExt(B',X)$ given by pulling back along $f$.

In a semi-abelian category $\C$, one can also define the notion internal action~\cite{BJK}. 
Internal actions on an object $X$ give rise to a functor
\[
\Act(-,X)\colon \C^{\op} \to \Set,
\]
and it was proved in~\cite{BJK} that there is a natural isomorphism
\[
\Act(-,X)\cong \SplExt(-,X).
\]
We do not describe internal actions in detail here, since split extensions provide a more practical framework, especially in the context of non-associative algebras.

\begin{definition}\cite{BJK}
A semi-abelian category $\C$ is said to be \emph{action representable} if the functor $\SplExt(-,X)$ is representable for every object $X$ of $\C$. That is, for every object $X$ there exists an object $[X]$ of $\C$ along with a natural isomorphism of functors
\[
\SplExt(-,X) \cong \Hom_{\C}(-,[X]).
\]
\end{definition}

Examples of action representable categories are the category $\Grp$ of groups and the category $\Lie$ of Lie algebras over a unital commutative ring. In the case of groups, every action of $B$ on $X$ corresponds to a group homomorphism $B \to \Aut(X)$, where $\Aut(X)$ is the automorphism group of $X$. Similarly, for Lie algebras, every action of~$B$ on $X$ is described by a Lie algebra homomorphism $B \to  \Der(X)$, where $\Der(X)$ is the Lie algebra of derivations on $X$.

However, action representability is a rather restrictive property. For instance, it was shown in~\cite[Theorem 6.1]{Tim} that amongst varieties of non-associative algebras over an infinite field $\F$ with $\bchar(\F) \neq 2$, the only action representable varieties are the category of abelian algebras and that of Lie algebras.

\begin{remark}\label{BJ}
It is shown in~\cite{BJK} that action representability is equivalent to the condition that, for every object $X$ in $\C$, the category $\SplExtX$ of split extensions with kernel $X$ has a terminal object of the form
\[
\begin{tikzcd}
X \arrow[r]
& {[X] \ltimes X} \arrow[r, shift left] &
{[X]} \arrow[l, shift left]
\end{tikzcd}
\]
\end{remark}

We can weaken this condition by assuming instead that every object in $\SplExtX$ is \emph{accessible}, i.e., it has a morphism into a \emph{subterminal} or \emph{faithful} object~\cite{act_accessible}, that is, an object which admits at most one morphism into it. This leads to the following notion, introduced by D.~Bourn and G.~Janelidze in~\cite{act_accessible} in order to compute centralizers of normal subobjects or of equivalence relations.

\begin{definition}\cite{act_accessible}
A semi-abelian category $\C$ is \emph{action accessible} if, for any object $X$ of $\C$, every object in $\SplExtX$ is accessible.
\end{definition}

The notion of action accessibility allows us to encompass a wider class of categories that are not action representable, such as the categories $\Rng$ and $\CRng$, the varieties $\Assoc$ and $\CAssoc$, and the variety $\Pois$ of Poisson algebras (see \Cref{SecPois}).

\begin{remark}
Since the existence of a terminal object in $\SplExtX$ is stronger than every object being accessible, it follows from \Cref{BJ} that
\[
\emph{action representability}\Rightarrow\emph{action accessibility.}
\]
\end{remark}

Alternatively, instead of weakening the condition on the existence of a terminal object in $\SplExtX$, one may weaken the properties of the representing object $[X]$. This is the approach taken in~\cite{Casas}, where the authors showed that every \emph{Orzech category of interest}~\cite{Orzech} admits a \emph{universal strict general actor} (USGA for short). A.~Montoli later proved in~\cite{Montoli} that all Orzech categories of interest are action accessible.

Recently in~\cite{WRA}, G.~Janelidze introduced the notion of \emph{weakly representable action}.

\begin{definition}\cite{WRA}
A semi-abelian category $\C$ is said to be \emph{weakly action representable} if the functor $\SplExt(-,X)$ admits a weak representation for every object~$X$ of~$\C$. This means that for every object $X$, there exists an object $T=T(X)$ of $\C$ and a natural monomorphism of functors
\[
\mu \colon \SplExt(-,X) \rightarrowtail \Hom_{\C}(-,T).
\]
A morphism $\varphi\colon B \to  T$ that belongs to $\Imm(\mu_B)$ is called \emph{acting morphism}.
\end{definition}

Examples of weakly action representable categories include the variety $\Assoc$ of associative algebras (see~\cite{WRA}), where $T=\Bim(X)$ is the associative algebra of \emph{bimultipliers} of $X$~\cite{MacLane58}; the variety $\Leib$ of Leibniz algebras (see~\cite[Theorem~3.14]{CigoliManciniMetere}), where $T=\Bider(X)$ is the Leibniz algebra of \emph{biderivations} of $X$~\cite{loday}; the variety $\Nil_2(\Com)$ and $\Nil_2(\ACom)$ of $2$-nilpotent commutative and anti-commutative algebras (see Theorems 2.19 and 2.21 of~\cite{WRAAlg}); and the variety $\CAssoc$ of commutative associative algebras (see~\cite[Theorem 2.11]{WRAAlg}). More in general, if in an Orzech category of interest $\C$ each $\USGA(X)$ is an object of the category, then $\C$ is weakly action
representable (see~\cite[Corollary 4.2]{CigoliManciniMetere}).

A result obtained by G.~Janelidze in~\cite{WRA} is that every weakly action representable category is action accessible. Thus, we have
\[
\textit{action representability}\Rightarrow\textit{weak action representability}\Rightarrow\textit{action accessibility.}
\]
The converse of the first implication does not hold, since the variety $\Assoc$ of associative algebras is not action representable (see~\cite[Proposition 1.11]{Tim}). J.~R.~A.~Gray recently observed in~\cite{Gray} that the converse of the second implication also fails: he proved that the varieties of $k$-solvable groups ($k \geq 3$) are not weakly action representable but action accessible. This result was then extended in~\cite{XabiMancini} to the varieties of $2$-solvable groups and $n$-nilpotent groups ($n \geq 3$). There, the authors also proved that the varieties of $k$-solvable Lie algebras ($k \geq 2$) and $n$-nilpotent Lie algebras ($n \geq 3$) fail to be weakly action representable.

\subsection*{External weak actor in varieties of non-associative algebras}

The representability of actions of non-associative algebras was studied in~\cite{WRAAlg}, where the authors proved that, for any object $X$ of an operadic, action accessible variety $\V$, there exists a \emph{partial algebra} $\E(X) \subseteq \End(X) \times \End(X)$, called the \emph{external weak actor} of $X$, together with a monomorphism of functors
\[
\tau \colon \SplExt(-,X) \rightarrowtail \Hom_\PAlg(\widetilde{U}(-),\E(X)),
\]
called \emph{external weak representation}. Here, $\PAlg$ denotes the category of \emph{partial algebras}, and $\widetilde{U} \colon \V \to \PAlg$ is the forgetful functor.

We recall the construction of $\E(X)$, since it will be central in proving that actions on a unital algebra $X$ are representable. We begin by recalling a result that links action accessibility with the identities of degree $3$ satisfied by a variety $\V$.

\begin{theorem}\cite{Tim, Coco_Thesis}\label{Theorem AC iff Orzech}
Let $\V$ be a homogeneous variety of non-associative algebras. The following conditions are equivalent:
\begin{tfae}
\item $\V$ is an Orzech category of interest;
\item $\V$ is action accessible;
\item There exist $\lambda_{1}, \dots, \lambda_{8}, \mu_{1}, \dots, \mu_{8} \in \F$ such that
\begin{align}
\begin{split}\label{eq00}
x(yz)=\lambda_1(xy)z & +\lambda_2(yx)z+\lambda_3 z(xy) + \lambda_4 z(yx)\\
& + \lambda_5 (xz)y + \lambda_6 (zx)y + \lambda_7 y(xz) + \lambda_8 y(zx)
\end{split}
\intertext{and}
\begin{split}\label{eq01}
(yz)x=\mu_1(xy)z & +\mu_2(yx)z+\mu_3 z(xy) + \mu_4 z(yx)\\
& + \mu_5 (xz)y + \mu_6 (zx)y + \mu_7 y(xz) + \mu_8 y(zx)
\end{split}
\end{align}
are identities of $\V$. The two previous identities together are called the \emph{$\lambda/\mu$-rules}. \noproof
\end{tfae}
\end{theorem}

\begin{remark}
Note that the choice of the $\lambda/\mu$-rules is not unique, but fixed for our purposes.
\end{remark}

Now, let $\V$ be an action accessible variety  determined by a set of multilinear identities
\[
\Phi_{k,i}(x_1,\ldots,x_k)=0, \quad i=1,\ldots,n,
\]
where $k$ is the degree of the polynomial $\Phi_{k,i}$, and fix $\lambda_1, \ldots,\lambda_8,\mu_1, \ldots,\mu_8 \in \F$ which determine a choice of the $\lambda/\mu$-rules. The external weak actor $\E(X)$ is defined as the subspace of all pairs $f=(f\ast-,-\ast f) \in \End(X) \times \End(X)$ satisfying
\[
\Phi_{k,i}(\alpha_1,\ldots,\alpha_k)=0, \quad \forall i=1,\ldots,n,
\]
for every choice of $\alpha_j=f$ and $\alpha_t \in X$, with $t\neq j\in \{1,\ldots,k\}$, where we write $fx\coloneq f\ast x$ and $xf\coloneq x\ast f$. 

Furthermore, $\E(X)$ is endowed with the \emph{bilinear partial operation}
\[
\langle -,- \rangle \colon \Omega \to \E(X)  \colon (f,g) \mapsto \langle f,g \rangle=h,
\]
where $\Omega=\langle -,- \rangle^{-1}(\E(X))$ is the preimage of the inclusion $\E(X) \hookrightarrow \End(X) \times \End(X)$, and $h=(h \ast -, - \ast h)$ is defined by
\begin{align*}
x \ast h=\lambda_1(x\ast f)\ast g & +\lambda_2(f\ast x)\ast g+\lambda_3g\ast(x\ast f) + \lambda_4 g \ast(f\ast x) \\
 & + \lambda_5 (x\ast g)\ast f + \lambda_6 (g\ast x)\ast f + \lambda_7 f\ast (x\ast g) + \lambda_8 f\ast (g\ast x)
\intertext{and}
h \ast x=\mu_1(x\ast f)\ast g     & +\mu_2(f\ast x)\ast g+\mu_3g\ast(x\ast f) + \mu_4 g\ast(f\ast x) \\
 & + \mu_5 (x\ast g)\ast f + \mu_6 (g\ast x)\ast f + \mu_7 f\ast (x\ast g) + \mu_8 f\ast (g\ast x).
\end{align*}
Finally, $\E(X)$ is called \emph{external actor} of $X$ whenever $\tau$ is a natural isomorphism.

\begin{example}\label{Example associative}
If $\V=\Assoc$ and we fix the standard choice of the $\lambda/\mu$-rules (that is $\lambda_1 = \mu_8 = 1$ and $\lambda_i=\mu_j=0$, for any $i \neq 1$ and $j \neq 8$), then $\E(X)$ is isomorphic to the associative algebra
\begin{align*}
&\Bim(X)
  = \{ (f \ast -,- \ast f) \in \End(X)\times \End(X)^{\text{op}} \mid \cdots \\ 
  &\quad \cdots \mid \; f \ast (xy)=(f \ast x)y,\,
  (xy) \ast f=x(y \ast f),\, x(f \ast y)=(x \ast f)y, \; \forall x,y \in X \}
\end{align*}
of \emph{bimultipliers} of $X$ (see~\cite{MacLane58}, where they are called \emph{bimultiplications}), where the multiplication is induced by the usual composition of functions in $\End(X)$. It was proved in~\cite{WRA} that $\Bim(X)$ provides a weak representation of the actions on $X$.
\end{example}

\begin{example}
If $\V=\CAssoc$, it was proved in~\cite[proof of Theorem 2.6]{BJK2} that there is a natural isomorphism
\[
\SplExt(-,X) \cong \Hom_\Assoc\left(\overline{U}(-),\M(X)\right),
\]
where $\overline{U} \colon \CAssoc \to \Assoc$ denotes the forgetful functor, and
\[
\M(X)=\{ f \in \End(X) \mid f(xy)=f(x)y, \; \forall x,y \in X \}
\]
is the associative algebra of \emph{multipliers} of $X$. Thus, $\E(X) \cong \M(X)$ is an external actor of $X$. We observe that, although $\M(X)$ is not in general a commutative algebra, it was shown in~\cite[Theorem 2.11]{WRAAlg} that the variety $\CAssoc$ is weakly action representable.
\end{example}

\begin{example}
Let $\bchar(\F) \neq 2$. If $\V=\Lie$, then $\E(X) \cong \Der(X)$ is the actor of~$X$.
\end{example}

\begin{remark}\label{rem_int}
For any algebra $X$ in $\V$, one defines a partial algebra homomorphism
\[
\Inn \colon \widetilde{U}(X) \to \E(X) \colon x \mapsto (L_x,R_x),
\]
where $L_x = x \cdot -$ and $R_x = - \cdot x$ denote, respectively, the left and right multiplications by $x$. The kernel of $\Inn$ is the \emph{annihilator}
\[
\Ann(X)=\{ x \in X \mid xy = yx = 0,\ \forall y \in X \}
\]
of $X$. In some cases, such as for Lie algebras, $\Ann(X)$ is called the \emph{center} of $X$ and is denoted by $\bZ(X)$.
\end{remark}

The following result explains how the construction of the external weak actor allows to represent actions on, or equivalently split extensions by, an algebra $X$.

\begin{theorem}\cite{WRAAlg}\label{thmV}
Let $\V$ be an operadic, action accessible variety of non-associative algebras, and fix a choice of the $\lambda/\mu$-rules.
\begin{enumerate}
\item Let $X$ be an object of $\V$. There exists a natural monomorphism of functors
\[
\tau \colon \SplExt(-,X) \rightarrowtail \Hom_{\PAlg}(\widetilde{U}(-),\E(X)),
\]
where $\widetilde{U} \colon \V \to \PAlg$ denotes the forgetful functor and, for every object $B$ of $\V$, $\tau_B$ is the injection which sends a split extension
\[
\begin{tikzcd}
X\arrow [r, "k"]
&A \arrow[r, shift left, "\alpha"] &
B\ar[l, shift left, "\beta"]
\end{tikzcd}
\]
to the partial algebra homomorphism
\[
\widetilde{U}(B) \to \E(X) \colon b \mapsto (b \ast -, - \ast b),
\]
where $b \ast x \coloneqq \beta(b)k(x)$ and $x \ast b \coloneqq k(x)\beta(b)$.
\item Let $B, X$ be objects of $\V$. The partial algebra homomorphism
\[
\widetilde{U}(B) \to \E(X) \colon b \mapsto (b \ast -,- \ast b)
\]
belongs to $\Imm(\tau_B)$ if and only if $\Phi_{k,i}(\alpha_1,\ldots,\alpha_k)=0$, where at least one of the $\alpha_1, \ldots, \alpha_k$ is an element $x \in X$, and the others are elements $b \in B$. Here, we write $bx\coloneq b\ast x$ and $xb\coloneq x\ast b$. 
\item If $\E(X)$ is an object of $\V$, then $(\E(X),\tau)$ becomes a weak representation of $\SplExt(-,X)$.
\noproof
\end{enumerate}
\end{theorem}

\begin{example}
Let $X$ be an object of $\V$, and consider the canonical action of $X$ on itself given by left and right multiplication, that is
\[
b \ast x = bx, \quad \text{and} \quad x \ast b = xb.
\]
Then the component $\tau_X$ sends this action to the partial algebra homomorphism $\Inn \colon \widetilde{U}(X) \to \E(X)$ of \Cref{rem_int}.
\end{example}

We conclude this section with the following proposition, showing that there are special cases where $\E(X)$ becomes the actor of $X$.

\begin{proposition}\label{prop_e(x)}
Let $\V$ be an operadic, action accessible variety of non-associative algebras, and fix a choice of $\lambda/\mu$-rules. For any object $X \in \V$, if the partial algebra homomorphism $\Inn \colon \widetilde{U}(X) \to \E(X)$ is an isomorphism, then the actions on $X$ are representable and $X$ is its own actor.
\end{proposition}

\begin{proof}
Suppose that $\Inn$ is an isomorphism. Then
\[
\Hom_{\PAlg}(\widetilde{U}(-),\E(X))=\Hom_{\V}(-,\E(X)) \cong \Hom_{\V}(-,X)
\]
and hence, by composition, we get a natural monomorphism of functors
\[
\widetilde{\tau} \colon \SplExt(-,X) \rightarrowtail \Hom_{\V}(-,X).
\]
It remains to show that, for any algebra $B$ of $\V$, the component $\widetilde{\tau}_B$ is surjective, that is, every homomorphism $\varphi \colon B \to X$ in $\V$ arises from a split extension of $B$ by~$X$ in $\V$.

To this end, consider the split extension
\begin{equation}\label{split_phi}
\begin{tikzcd}
X \arrow[r, "\iota_2"] &
B \ltimes X \arrow[r, shift left, "\pi_1"] &
B, \arrow[l, shift left, "\iota_1"]
\end{tikzcd}
\end{equation}
where $B \ltimes X$ is the algebra whose underlying vector space is $B \times X$ and whose multiplication is defined by
\[
(b,x) \cdot (b',x') = (bb', xx' + \varphi(b)x' + x\varphi(b')),
\]
and $\iota_1, \iota_2$ and $\pi_1$ are the canonical injections and projection. Straightforward computations show that $\tau_B$ sends \eqref{split_phi} to the morphism
\[
\Inn \circ \varphi \colon B \to \E(X) \colon b \mapsto (L_{\varphi(b)}, R_{\varphi(b)}).
\]
Consequently, $\widetilde{\tau}_B$ sends the split extension \eqref{split_phi} to the morphism $\varphi$.

Finally, we need to check that \eqref{split_phi} is indeed a split extension in $\V$. Let $\Phi(x_1,\ldots,x_k)=0$ be any multilinear identity of $\V$. Then, for any $\alpha_1,\ldots,\alpha_k \in B \cup X$, the evaluation $\Phi(\alpha_1,\ldots,\alpha_k)$ reduces to $0$, since
\[
b \ast x = \varphi(b)x \; \text{ and } \; x \ast b = x\varphi(b),
\]
for every $b \in B$ and $x \in X$. Thus, any such evaluation reduces to an expression involving only multiplications in $X$, and hence vanishes. It follows that $B \ltimes X$ satisfies all identities of $\V$, and therefore is an algebra of $\V$.

Therefore, $\widetilde{\tau}_B$ is surjective, and the actions on $X$ are representable. Since $\Inn$ is an isomorphism, it follows that $X$ is its own actor.
\end{proof}

\begin{example}\label{ex_assoc}
Let $\V=\Assoc$ and fix the standard choice of the $\lambda/\mu$-rules. Then, for any associative algebra $X$, the map
\[
\Inn \colon X \to \Bim(X) \colon x \mapsto (L_x,R_x)
\]
is an isomorphism if and only if $X$ is a unital algebra.

Indeed, if $\Inn$ is an isomorphism, then $X$ is a unital algebra with unit
\[
1=\Inn^{-1}(\id_X,\id_X).
\]
Conversely, if $X$ has unit $1$, then $\Ker(\Inn)=\Ann(X)=\{ 0 \}$, and $\Inn$ is injective. Furthermore, as observed in \cite[Remark 3.8]{WRAAlg}, any bimultiplier of $X$ is uniquely determined by the choice of an element $x=f \ast 1 = 1 \ast f \in X$, since
\[
f \ast y = xy \quad \text{and} \quad y \ast f = yx
\]
for any $y \in X$. Thus, $\Inn$ is surjective and, by \Cref{prop_e(x)}, there is a natural isomorphism
\[
\SplExt(-,X) \cong \Hom_\Assoc(-,X).
\]
Hence, actions on a unital associative algebra $X$ are representable. Similarly, one can show that the same result holds for the variety $\CAssoc$ of commutative associative algebras, replacing $\Bim(X)$ with the associative algebra $\M(X)$ of multipliers of $X$.
\end{example}

\begin{example}\label{ex_jjord}
Let $\V=\JJord$ be the variety of Jacobi-Jordan algebras over a field of characteristic $3$ and let $X$ be an object of $\V$. We recall that, in this case, the external actor $\E(X)$ is isomorphic to the partial algebra
\[
\ADer(X)=\{ f \in \End(X) \mid f(xy) = -f(x)y - f(y)x, \; \forall x,y \in X \},
\]
of \emph{anti-derivations} of $X$, endowed with the bilinear partial operation
\[
\{ f,g \} = -f \circ g - g \circ f,
\]
called the \emph{anti-commutator} of $f$ and $g$. Again, we prove that the map
\[
\Inn \colon \widetilde{U}(X) \to \ADer(X) \colon x \mapsto L_x
\]
is an isomorphism if and only if $X$ is a unital Jacobi-Jordan algebra.

Indeed, if $\Inn$ is an isomorphism, then $X$ is a unital algebra with unit
\[
1=\Inn^{-1}(\id_X).
\]
Conversely, if $X$ has unit $1$, then $\Ker(\Inn)=\Ann(X)=\{ 0 \}$, and $\Inn$ is injective. Furthermore, any anti-derivation of $X$ is uniquely determined by the choice of an element $x=f(1)$, since
\[
f(y)=f(y \cdot 1)=-f(y)1-f(1)y
\]
for any $y \in X$, and hence
\[
f(y)=f(1)y=xy.
\]
Thus, $\Inn$ is surjective and, by \Cref{prop_e(x)}, the actions on a unital Jacobi-Jordan algebra $X$ are representable.
\end{example}

In \Cref{SecMain}, we will show that the above proposition applies to all unital algebras in any operadic, action accessible, unit-closed variety $\V$.

\section{Representability of actions in ideally exact categories}\label{SecIdEx}

Let $\V$ be a unit-closed variety. In order to study the representability of actions of unital non-associative algebras, we need to recall the definition of an \emph{ideally exact category}, introduced by G.~Janelidze in~\cite{IdeallyExact}.
 
\begin{definition}\cite{IdeallyExact}
Let $\U$ be a category with pullbacks, with initial object $0$ and terminal object $1$. $\U$ is said to be \textit{ideally exact} if it is Barr-exact, Bourn-protomodular, has finite coproducts and the unique morphism $0 \to 1$ in $\U$ is a regular epimorphism.
\end{definition}

In addition to all semi-abelian categories, examples of ideally exact categories include the categories $\Ring$ and $\CRing$ of unital rings and unital commutative rings, all non-trivial protomodular varieties of universal algebras, such as the variety of \emph{MV-algebras}~\cite{rel}, and any cotopos. 

Moreover, the following characterisation holds.
 
\begin{theorem}\cite{IdeallyExact}\label{carid}
Let $\U$ be a category with pullbacks. The following conditions are equivalent:
\begin{enumerate}
\item $\U$ is ideally exact.
\item $\U$ is Barr-exact, has finite coproducts and there exists a monadic functor $\U \to \V$, where $\V$ is a semi-abelian category.
\item There exists a monadic functor $\U \to \V$, where $\V$ is a semi-abelian category, such that the underlying functor of the corresponding monad preserves regular epimorphisms and kernel pairs. \noproof
\end{enumerate}
\end{theorem}

\begin{remark}\cite{IdeallyExact}\label{cartesian}
Let $\U$ be an ideally exact category. It follows from \Cref{carid} that there exists a monadic adjunction 
\begin{equation}\label{monadic}
\begin{tikzcd}
{\U} & {\V}
\arrow[""{name=0, anchor=center, inner sep=0}, "U"', from=1-1, to=1-2]
\arrow[""{name=1, anchor=center, inner sep=0}, "F"', curve={height=12pt}, from=1-2, to=1-1]
\arrow["\dashv"{anchor=center, rotate=-90}, draw=none, from=1, to=0]
\end{tikzcd}
\end{equation}
with $\V$ semi-abelian. By~\cite[Theorems 2.4 and 2.6]{IdeallyExact}, this adjunction is associated with the unique morphism $0 \to 1$ (up to an equivalence) if and only if the unit of the adjunction is cartesian. 

Of course the monadic adjunction associated to $0 \to 1$ is always there. Indeed, one may take $\V = (\U \downarrow 0)$, the slice category over $0$, where $U=!^*$ is the pullback functor along $0 \to 1$, and $F$ is the functor which sends a morphism $X \to 0$ to its domain. However, this construction is not always the most natural choice. In particular, if $\U$ is already semi-abelian, it may be preferable to take $\V=\U$. We further observe that, since $F$ is a left adjoint, $F(0)$ is an initial object of $\U$.
\end{remark}

\begin{example}\label{monadic_ring}
Let $\U=\Ring$, which has the ring of integers $\Z$ as initial object and the zero ring $\{0\}$ as terminal one. Then a monadic adjunction as in \Cref{cartesian} is given by
\begin{equation*}\label{eq:situation_rng}
\begin{tikzcd}
{\Ring} & {\Rng,}
\arrow[""{name=0, anchor=center, inner sep=0}, "U"', from=1-1, to=1-2]
\arrow[""{name=1, anchor=center, inner sep=0}, "F"', curve={height=12pt}, from=1-2, to=1-1]
\arrow["\dashv"{anchor=center, rotate=-90}, draw=none, from=1, to=0]
\end{tikzcd}
\end{equation*}
where $U$ is the forgetful functor and $F$ maps any (not necessarily unital) ring $X$ to the semidirect product~$\Z \ltimes X$ with operations
\[
(n,x)+(m,y)=(n+m,x+y), \quad (n,x) \cdot (m,y)=(nm,xy+ny+mx)
\]
and unit element $(1,0)$. The unit $\eta \colon 1_{\Rng} \Rightarrow UF$ of the adjunction is cartesian since the map $\eta_X \colon X \to  U(\Z \ltimes X) \colon x \mapsto (0,x)$ is the kernel of
\[
U(\pi_1) \colon U(\Z \ltimes X) \to  U(\Z) \colon (n,x) \mapsto n,
\]
Thus, $F \dashv U$ is, up to an equivalence, the adjunction associated with the unique morphism $\Z \to \{ 0 \}$ in $\Ring$.

Similarly, there is a monadic adjunction with cartesian unit
\[
\begin{tikzcd}
{\CRing} & {\CRng.}
\arrow[""{name=0, anchor=center, inner sep=0}, "U"', from=1-1, to=1-2]
\arrow[""{name=1, anchor=center, inner sep=0}, "F"', curve={height=12pt}, from=1-2, to=1-1]
\arrow["\dashv"{anchor=center, rotate=-90}, draw=none, from=1, to=0]
\end{tikzcd}
\]
\end{example}

\begin{example}\label{monadic_alg}
Let $\V$ be a unit-closed variety. The class $\V_1$ of unital algebras of~$\V$ can be seen as a category whose morphisms are the algebra homomorphisms preserving the unit. Then $\V_1$ is a protomodular variety of universal algebras. Hence, by~\cite[Example 3.6]{IdeallyExact}, it is an ideally exact category. Moreover, $\V_1$ is not pointed, since an initial object is the field $\F$ endowed with the standard multiplication, while the terminal one is the zero algebra $\{0\}$.

Furthermore, there exists a monadic adjunction with cartesian unit.
\[
\begin{tikzcd}
{\V_1} & {\V,}
\arrow[""{name=0, anchor=center, inner sep=0}, "U"', from=1-1, to=1-2]
\arrow[""{name=1, anchor=center, inner sep=0}, "F"', curve={height=12pt}, from=1-2, to=1-1]
\arrow["\dashv"{anchor=center, rotate=-90}, draw=none, from=1, to=0]
\end{tikzcd}
\]
where $U$ is the forgetful functor and $F$ maps any algebra $X$ of $\V$ to the semidirect product $\F \ltimes X$, with operations
\[
(n,x)+(m,y) = (n+m, x+y), \quad (n,x)\cdot(m,y) = (nm, xy + ny + mx)
\]
and unit $(1,0)$. 
\end{example}

In~\cite{IdeallyExact}, G.~Janelidze redefined the notion of action representability in the setting of ideally exact categories.

\begin{definition}\cite{IdeallyExact}\label{def_act_repr}
Let $\U$ be an ideally exact category and let $!^* \colon \U \to (\U \downarrow 0)$ be the pullback functor along $0 \to 1$. Then $\U$ is \emph{action representable} if the functor
\[
\SplExt(-,!^*(X)) \cong \Act(-,!^*(X)) \colon (\U \downarrow 0)^{\op} \to \Set
\]
is representable for any object $X$ of $\U$.
\end{definition}

\begin{example}\cite{IdeallyExact}\label{ex_ring}
Let $\V=\Rng$ be the (semi-abelian) category of not necessarily unital rings. As it is shown in~\cite{BJK2}, actions on a ring $X$ are representable, in particular, when $X$ is unital. In this case, the canonical ring homomorphism
\[
X \to \Bim(X) \colon x \mapsto (L_x,R_x),
\]
where $\Bim(X)$ is the ring of \emph{bimultipliers} of $X$~\cite{MacLane58}, is an isomorphism. 

Now, recall there is a canonical category equivalence $(\Ring \downarrow \Z) \simeq \Rng$, which is defined as follows: for any object $X \to \Z$ of $(\Ring \downarrow \Z)$, its kernel is an object of~$\Rng$; conversely, every not necessarily unital ring $X$ can be seen as the kernel of the projection $\Z \ltimes X \to \Z \colon (n,x) \mapsto n$.

Under this equivalence, the pullback functor $!^* \colon \Ring \to (\Ring \downarrow \Z)$ along the unique morphism $\Z \to \{ 0 \}$ is the same as the forgetful functor $U \colon \Ring \to \Rng$. Indeed, given any unital ring $X$, viewed as the object $X \to \{ 0 \}$ of $(\Ring \downarrow \{ 0 \})$, its image under $!^*$ is given by the pullback
\[
\begin{tikzcd}
\Z \times X \arrow[r, "\pi_2"] \arrow[d, swap, "\pi_1"] \arrow[dr, phantom, "\usebox\pullback" , very near start, color=black] & X \arrow[d] \\
\Z \arrow[r] & \{ 0 \}
\end{tikzcd}
\]
and one has that $\Ker \pi_1 \cong X$, seen as a ring without unit. Since, under the equivalence $(\Ring \downarrow \Z) \simeq \Rng$, any object $A \to \Z$ corresponds to its kernel, it follows that $!^*$ acts by forgetting the unit.

Thus, given any unital ring $X$, the functor $\SplExt(-, !^{*}(X)) \colon (\Ring \downarrow \Z)^{\op} \to \Set$ is representable if and only if so is functor $\SplExt(-,U(X)) \colon \Rng^{\op} \to  \Set$. Hence, as observed in~\cite[Section 5]{IdeallyExact}, we conclude that $\Ring$ is action representable, and that the actor of a unital ring $X$ is $X$ itself. The same argument applies to the category $\CRing$ of unital commutative rings.
\end{example}

\Cref{ex_ring} suggests that, in order to study the representability of actions of the category of unital rings, it is more convenient to choose the monadic adjunction with cartesian unit of \Cref{monadic_ring}.  In a similar way, for any unit-closed variety of non-associative $\F$-algebras $\V$, there is a category equivalence $(\V_1 \downarrow \F) \simeq \V$. Thus, actions on a unital algebra $X$ are representable if and only if the functor
\[
\SplExt(-,U(X)) \colon \V^{\op} \to \Set,
\]
where $U \colon \V_1 \to \V$ denotes the forgetful functor, is representable.

\begin{remark}
It is straightforward to verify that, if a variety $\V$ is action representable, then so is its subcategory $\V_1$ of unital algebras. This applies, for instance, to the category $\BoolAlg$ of \emph{Boolean algebras}, which can be described as the subcategory of Boolean rings with unit and unit-preserving morphisms (see~\cite{Bell}).

It was shown in~\cite[Proposition~3.1]{BJK2} that the variety of Boolean rings is action representable, and that the actor of a Boolean ring $X$ is the ring $\End(X)$ of $X$-linear endomorphisms of $X$. In particular, it follows that actions on a unital Boolean ring~$X$ are representable.

Moreover, in this case, the actor of a Boolean algebra $X$ can be identified with~$X$ itself. Indeed, any $X$-linear endomorphism $f \colon X \to X$ is uniquely determined by the element $x = f(1) \in X$, since for every $y \in X$ one has $f(y) = f(1)y = xy$.

Thus, the ideally exact category $\BoolAlg$ is action representable, and the actor of any Boolean algebra $X$ is $X$ itself.
\end{remark}

Following~\cite{IdeallyExact}, one may also re-define the notions of action accessibility and weak action representability.

\begin{definition}\label{act_acc}
Let $\U$ be an ideally exact category and let $!^* \colon \U \to (\U \downarrow 0)$ be the pullback functor along $0 \to 1$. Then
\begin{itemize}
\item $\U$ is \emph{action accessible} if, for any object $X$ of it, every object in the category $\SplExt(!^*(X))$ of split extensions with kernel $!^*(X)$ is accessible.
\item $\U$ is \emph{weakly action representable} if the functor $\SplExt(-,!^*(X))$ admits a weak representation for any object $X$ of $\U$.
\end{itemize}
\end{definition}

\begin{remark}
We observe that the chain of implications
\[
\textit{action representability}\Rightarrow\textit{weak action representability}\Rightarrow\textit{action accessibility}
\]
still holds in the setting of ideally exact categories. In fact:
\begin{itemize}
\item If $\U$ is action representable, then it is immediate that it is weakly action representable.
\item If $\U$ is weakly action representable, then Corollary~4.3 and Proposition~4.4 of~\cite{WRA} apply, in particular, to all split extensions in $(\U \downarrow 0)$ with kernel of the form $!^*(X)$, for some object $X$ of $\U$. Consequently, every object in $\SplExt(!^*(X))$ is accessible, and hence $\U$ is an action accessible category. We point the reader to~\cite[Section 2]{Mancini_Thesis}, where an explicit proof of this implication for semi-abelian categories is given.
\end{itemize}
\end{remark}

In the next section, we prove that, in the context of unital non-associative algebras, the existence of the $\lambda/\mu$-rules implies action representability. Our main tool is to show that the existence of a multiplicative unit is equivalent to the canonical map $\Inn \colon \widetilde{U}(X) \to \E(X)$ being an isomorphism.

Before doing so, we show that action accessibility fails in general in the setting of categories of unital algebras, as the following proposition illustrates. In this setting, as in the case of action representability, we replace the pullback functor along $\F \to \{0\}$ with the forgetful functor $U \colon \V_1 \to \V$. In other words, we work with the monadic adjunction with cartesian unit described in \Cref{monadic_alg}.

\begin{proposition}
The ideally exact categories $\Alg_1$ and $\Com_1$ are not action accessible.
\end{proposition}

\begin{proof}
Let $\V=\Com$ be the category of commutative (not necessarily associative) algebras, and consider the monadic adjunction of \Cref{monadic_alg}. We want to show there is an object $X$ of $\Com_1$ and a split extension
\[
\begin{tikzcd}
U(X) \arrow [r, "k"]
& A \arrow[r, shift left, "\alpha"] &
B \ar[l, shift left, "\beta"]
\end{tikzcd}    
\]
in $\Com$ which does not admit any morphism into a faithful object of $\SplExt(U(X))$. Our construction is inspired by the argument in~\cite[Proposition 6.4]{Cappelletti}, where it is shown that the category of \emph{lattice-ordered groups} is not action accessible. The idea is to adapt the technique used there to the setting of non-associative algebras.  

We take $X$ to be the $6$-dimensional unital commutative algebra with basis $\{ 1,e_1,e_2,e_3,f_1,f_2 \}$, where $1$ is the unit of $X$, and whose multiplication table is defined by
\[
e_i^2=f_j^2=1, \quad e_if_j=f_je_i=0, \quad f_1 f_2 = f_2 f_1=e_1
\]
and
\[
e_1 e_2 = e_2 e_1=e_3, \quad e_2e_3=e_3e_2=e_1, \quad e_3e_1=e_1e_3=e_2.
\]
We observe that $X$ is a simple algebra, i.e., its only ideals are $\{ 0 \}$ and $X$ itself. Indeed, let $I$ be a non-zero ideal of $X$, and let 
\[
0 \neq x=a1+\sum_{i=1}^3 b_i e_i + \sum_{j=1}^2 c_j f_j \in I,
\]
with $a,b_1,b_2,b_3,c_1,c_2 \in \F$. Then
\[
((xf_1)e_1)f_1=((af_1+c_1 1+c_2e_1)e_1)f_1=(c_1e_1+c_2 1)f_1=c_2f_1 \in I.
\]
If $c_2 \neq 0$, then $1=f_1^2 \in I$ and $I=X$. If instead $c_2=0$, then
\[
(x f_1)e_1=c_1 e_1 \in I.
\]
Again, if $c_1 \neq 0$, then $1=e_1^2 \in I$ and $I=X$. If instead $c_1=c_2=0$, then
\[
x f_1=a f_1 \in I.
\]
If $a \neq 0$, then $1 \in I$ and $I=X$. Now, suppose that $a=c_1=c_2=0$. In this case
\[
x=b_1 e_1 + b_2 e_2 + b_3 e_3
\]
with at least one $b_i \neq 0$. Without loss of generality, we suppose that $b_1 \neq 0$. Then 
\[
((xe_1)f_1)f_1=((b_1 1 + b_2 e_3 + b_3 e_2)f_1)f_1=b_1f_1^2=b_1 1\in I,
\]
which implies that $I=X$. Thus, $X$ is a simple algebra.

Now, consider the action of $U(X)$ on itself given by left and right multiplications, that is, $b \ast x=bx=xb=x \ast b$ for any $b,x \in X$. This gives rise to a split extension
\begin{equation}\label{non_faithful}
\begin{tikzcd}
U(X) \arrow [r, "\iota_2"]
& U(X) \ltimes U(X) \arrow[r, shift left, "\pi_1"] &
U(X) \ar[l, shift left, "\iota_1"]
\end{tikzcd}        
\end{equation}
in $\Com$, where $\iota_1,\iota_2$ and $\pi_1$ are the canonical injections and projection, and the bilinear multiplication on $U(X) \ltimes U(X)$ is defined by
\[
(b,x) \cdot (b',x')=(bb',xx'+bx'+xb').
\]
Now, we consider the algebra automorphism $\varphi \colon X \to X$ defined by
\[
\varphi(1)=1, \quad \varphi(e_1)=e_1, \quad \varphi(e_2)=-e_2, \quad \varphi(e_3)=-e_3, \quad \varphi(f_1)=f_1, \quad \varphi(f_2)=f_2.
\]
Thus, there are two morphisms in $\SplExt(U(X))$:
\[
\begin{tikzcd}
U(X) \arrow [r, "\iota_2"] \ar[d, equal]
&U(X) \ltimes U(X) \arrow[r, shift left, "\pi_1"] \ar[d, "\id_X \times \id_X"']&
U(X) \ar[l, shift left, "\iota_1"]\ar[d, "\id_X"] \\
U(X) \arrow [r, "\iota_2"]
& U(X) \ltimes U(X) \arrow[r, shift left, "\pi_1"] &
U(X) \ar[l, shift left, "\iota_1"]
\end{tikzcd}
\]
and
\[
\begin{tikzcd}
U(X) \arrow [r, "\iota_2"] \ar[d, equal]
&U(X) \ltimes U(X) \arrow[r, shift left, "\pi_1"] \ar[d, "\varphi \times \id_X"']&
U(X) \ar[l, shift left, "\iota_1"]\ar[d, "\varphi"] \\
U(X) \arrow [r, "\iota_2"]
& U(X) \ltimes U(X) \arrow[r, shift left, "\pi_1"] &
U(X) \ar[l, shift left, "\iota_1"]
\end{tikzcd}
\]
and we deduce that \eqref{non_faithful} is not faithful.

So, if $\Com_1$ were action
accessible, then there should exist a faithful object
\[
\begin{tikzcd}
U(X) \arrow [r, "k"]
& A \arrow[r, shift left, "\alpha"] &
B \ar[l, shift left, "\beta"]
\end{tikzcd}
\]
in $\SplExt(U(X))$ and a morphism
\[
\begin{tikzcd}
U(X) \arrow [r, "\iota_2"] \ar[d, equal]
&U(X) \ltimes U(X) \arrow[r, shift left, "\pi_1"] \ar[d, "g"']&
U(X) \ar[l, shift left, "\iota_1"]\ar[d, "f"] \\
U(X) \arrow [r, "k"]
& A \arrow[r, shift left, "\alpha"] &
B \ar[l, shift left, "\beta"] .
\end{tikzcd}
\]
Then, if we consider the regular epimorphism-monomorphism factorization of $(g,f)$, we get a morphism of split extensions
\begin{equation}\label{splext_morph}
\begin{tikzcd}
U(X) \arrow [r, "\iota_2"] \ar[d, equal]
&U(X) \ltimes U(X) \arrow[r, shift left, "\pi_1"] \ar[d, twoheadrightarrow, "g"']&
U(X) \ar[l, shift left, "\iota_1"]\ar[d, twoheadrightarrow, "f"]\\
U(X) \arrow [r, "\widetilde k"]
& \Imm g \arrow[r, shift left, "\widetilde \alpha"] &
\Imm f \ar[l, shift left, "\widetilde \beta"],
\end{tikzcd}
\end{equation}
where $\Imm f \cong U(X)/\Ker f$. However, $X$ has only two ideals: $\{ 0 \}$ and $X$. Hence we have two possibilities: $\Imm f \cong U(X)$ or $\Imm f \cong \{ 0 \}$. If $\Imm f \cong U(X)$, then $f$ would be injective and so by~\cite[Proposition 1.4]{act_accessible} the split extension \eqref{non_faithful} would be faithful, and this is a contradiction. Alternatively, if $\Imm f \cong \{ 0 \}$, then $f$ is the zero morphism, and so $\Imm g \cong U(X)$. Therefore, recalling from~\cite{malcev} that the square of the right of \eqref{splext_morph} is a pullback, we get a contradiction since $U(X) \ltimes U(X)$ and the (direct) product $U(X) \times U(X)$ are not isomorphic $\F$-algebras. 

Thus, we conclude that $\Com_1$ is not action accessible. Furthermore, since the algebra $X$ is also an object of $\Alg_1$, it follows that the category $\Alg_1$ of all unital not necessarily associative algebras is not action accessible.
\end{proof}

\section{Action representability of unital algebras}\label{SecMain}

The aim of this section is to prove that, given an operadic, action accessible, unit-closed variety of non-associative algebras $\V$, the ideally exact category $\V_1$ is action representable. We start by proving the following proposition which establishes a connection between action accessibility and the property of being a unit-closed variety.

\begin{proposition}\label{lemma_unit}
Let $\V$ be an operadic, action accessible, unit-closed variety of non-associative algebras, and fix a choice of the $\lambda/\mu$-rules.
\begin{enumerate}
\item If $\V$ satisfies an identity of degree $2$, then 
\begin{itemize}
\item if $\bchar(\F)\neq 3$, $\V=\CAssoc$;
\item if $\bchar(\F)= 3$, $\V=\CAssoc$ or $\V$ is a subvariety of $\JJord$.
\end{itemize}
\item If $\V$ satisfies no identities of degree $2$, then 
\begin{equation}\label{eq_lemma1}
\begin{alignedat}{2}
\lambda_1 &= 1 + \lambda_4, \qquad & \lambda_2 &= -\lambda_4 - \lambda_7 - \lambda_8, \\
\lambda_3 &= -\lambda_4 - \lambda_6 - \lambda_8, \qquad & \lambda_5 &= \lambda_8
\end{alignedat}
\end{equation}
and
\begin{equation}\label{eq_lemma2}
\begin{alignedat}{2}
\mu_1 &= \mu_4, \qquad & \mu_2 &= 1 - \mu_4 - \mu_7 - \mu_8, \\
\mu_3 &= 1 - \mu_4 - \mu_6 - \mu_8, \qquad & \mu_5 &= -1 + \mu_8.
\end{alignedat}
\end{equation}
\item For any algebra $X$ of $\V$, the pair $(\id_X,\id_X)$ is the unit of the external weak actor $\E(X)$.
\end{enumerate}
\end{proposition}

\begin{proof}{\ }
\begin{enumerate}
\item Suppose that $\V$ satisfies an identity of degree $2$. Since $\V$ is operadic, such an identity must be of the form
\[
xy+\lambda yx=0.
\]
By~\cite[Proposition 1.3]{Tim}, one has $\lambda=0$ or $\lambda=\pm 1$. Since $\V$ is unit-closed, we must have $\lambda=-1$, and therefore $\V$ is a subvariety of $\Com$. Consequently, the $\lambda/\mu$-rules reduce to a single equation of the form
\[
x(yz)=\alpha(xy)z+\beta(xz)y
\]
for some $\alpha,\beta \in \F$. By~\cite[Proposition~2.1]{WRAAlg}, it follows that $\V$ is either a subvariety of $\CAssoc$, or a subvariety of $\JJord$. However, 
if $\bchar(\F)\neq 3$, the variety $\JJord$ is not unit-closed, and hence this case is excluded. Therefore, in this case, $\V$ must be a subvariety of $\CAssoc$. However, the only proper operadic subvarieties of $\CAssoc$ are the varieties of $k$-nilpotent commutative associative algebras, with $k \geq 1$,  
which are not unit-closed. Therefore, we conclude that $\V = \CAssoc$.

\item Suppose that $\V$ does not satisfy any identity of degree $2$, and let $X$ be the free algebra on two generators $a,b$. Since $\V$ is unit-closed, the unital algebra $\widetilde X=\langle X,1\rangle$ is again an object of $\V$. Applying the $\lambda/\mu$-rules in $\widetilde X$ with $(x,y,z)=(1,a,b)$, we obtain
\[
ab=(\lambda_1+\lambda_2+\lambda_7+\lambda_8)\,ab +(\lambda_3+\lambda_4+\lambda_5+\lambda_6)\,ba
\]
and
\[
ab=(\mu_1+\mu_2+\mu_7+\mu_8)\,ab +(\mu_3+\mu_4+\mu_5+\mu_6)\,ba.
\]
Since there are no identities of degree $2$, then $ab$ and $ba$
are linearly independent, and it follows that
\begin{equation}\label{system1}
\begin{alignedat}{2}
\lambda_1+\lambda_2+\lambda_7+\lambda_8 &= 1, \qquad &
\lambda_3+\lambda_4+\lambda_5+\lambda_6 &= 0, \\
\mu_1+\mu_2+\mu_7+\mu_8 &= 1, \qquad &
\mu_3+\mu_4+\mu_5+\mu_6 &= 0.
\end{alignedat}
\end{equation}
Similarly, by applying the $\lambda/\mu$-rules respectively with
$(x,y,z)=(a,1,b)$ and $(x,y,z)=(a,b,1)$, one gets
\begin{equation}\label{system2}
\begin{alignedat}{2}
\lambda_1+\lambda_2+\lambda_5+\lambda_7 &= 1, \qquad &
\lambda_3+\lambda_4+\lambda_6+\lambda_8 &= 0, \\
\mu_3+\mu_4+\mu_6+\mu_8 &= 1, \qquad &
\mu_1+\mu_2+\mu_5+\mu_7 &= 0.
\end{alignedat}
\end{equation}
and
\begin{equation}\label{system3}
\begin{alignedat}{2}
\lambda_1+\lambda_3+\lambda_5+\lambda_6 &= 1, \qquad &
\lambda_2+\lambda_4+\lambda_7+\lambda_8 &= 0, \\
\mu_2+\mu_4+\mu_7+\mu_8 &= 1, \qquad &
\mu_1+\mu_3+\mu_5+\mu_6 &= 0.
\end{alignedat}
\end{equation}
Solving the system given by equations
\eqref{system1}--\eqref{system2}--\eqref{system3}, one obtains
\eqref{eq_lemma1} and \eqref{eq_lemma2}.


\item Let $X$ be a not necessarily unital algebra of $\V$. If $\V$ satisfies an identity of degree $2$, then $\V=\CAssoc$ when $\bchar(\F)\neq 3$, while for $\bchar(\F)=3$ one may also have $\V \subseteq \JJord$. If $\V=\CAssoc$, by~\cite[Theorem~3.9(5)]{WRAAlg} the linear map
\[
\varphi \colon f \mapsto (f,f)
\]
defines an isomorphism $\M(X) \cong \E(X)$. Since the identity map $\id_X$ is the unit of the algebra of multipliers $\M(X)$, it follows that $\varphi(\id_X)=(\id_X,\id_X)$ is the unit of $\E(X)$. 

In a similar way, if $\bchar(\F)=3$ and $\V$ is a subvariety of $\JJord$, then~$\E(X)$ is isomorphic to a subalgebra of the partial algebra $\ADer(X)$ (see~\Cref{ex_jjord}). One has that $\id_X=\Inn(1) \in \ADer(X)$. Furthermore, $\id_X$ is the unit of $\ADer(X)$ since
\[
\langle f,\id_X \rangle=-f \circ \id_X - \id_X \circ f=-2f=f=\langle \id_X,f \rangle,
\]
for any $f \in \ADer(X)$.

Suppose now that $\V$ does not satisfy any identity of degree $2$. We first claim that $(\id_X,\id_X)\in\E(X)$. If $X$ is unital, then
\[
(\id_X,\id_X) = (L_{1_X},R_{1_X}) = \Inn(1_X) \in \E(X).
\]
If $X$ does not have a multiplicative unit, we consider the unital algebra $\widetilde X=\langle X,1 \rangle$. Then
\[
(\id_{\widetilde X},\id_{\widetilde X}) = (L_1,R_1) = \Inn(1) \in \E(\widetilde X).
\]
By definition of $\E(X)$, this means that for every identity $\Phi(x_1,\ldots,x_k)=0$ which determines $\V$ we have
\begin{equation}\label{eq_identity}
\Phi(\alpha_1,\ldots,\alpha_k)=0,
\end{equation}
for any choice of $\alpha_j=\id_{\widetilde X}$ and $\alpha_t\in\widetilde X$, with $t \neq j \in 1,\ldots,k$. Since $X \subseteq \widetilde X$, equation~\eqref{eq_identity} holds in particular for all $\alpha_t\in X$. Therefore $(\id_X,\id_X)$ satisfies all the identities of $\V$, that is, $(\id_X,\id_X) \in \E(X)$.

Finally, recall that $(\id_X,\id_X)$ is the unit of $\E(X)$ if and only if
\[
\langle f, (\id_X, \id_X) \rangle = \langle (\id_X, \id_X), f \rangle = f,
\]
for any $f = (f \ast -, - \ast f) \in \E(X)$. Equivalently, this condition is expressed by
\begin{align*}
x \ast f
&= (\lambda_1+\lambda_3+\lambda_5+\lambda_6)(x \ast f)
 + (\lambda_2+\lambda_4+\lambda_7+\lambda_8)(f \ast x)\\
&= (\lambda_1+\lambda_2+\lambda_5+\lambda_7)(x \ast f)
 + (\lambda_3+\lambda_4+\lambda_6+\lambda_8)(f \ast x)
\end{align*}
and
\begin{align*}
f \ast x
&= (\mu_1+\mu_3+\mu_5+\mu_6)(x \ast f)
 + (\mu_2+\mu_4+\mu_7+\mu_8)(f \ast x)\\
&= (\mu_1+\mu_2+\mu_5+\mu_7)(x \ast f)
 + (\mu_3+\mu_4+\mu_6+\mu_8)(f \ast x)
\end{align*}
for any $x \in X$. One may check that these equalities are direct consequences of the system of equations~\eqref{eq_lemma1} and \eqref{eq_lemma2}.
\end{enumerate}
\end{proof}

\begin{example}
Let $\V=\Assoc$ and fix the standard choice of the $\lambda/\mu$-rules as in \Cref{Example associative}. Then, one may check that \eqref{eq_lemma1}--\eqref{eq_lemma2} are satisfied.
\end{example}

\begin{example}\label{example Alt}
Let $\bchar(\F) \neq 2$ and consider the variety $\Alt$ of alternative algebras. The multilinearisation process shows that $\Alt$ is equivalent to the variety determined by
\[
(xy)z+(xz)y-x(yz)-x(zy)=0
\]
and
\[
(xy)z+(yx)z-x(yz)-y(xz)=0.
\] 
A choice of the $\lambda/\mu$-rules is then given by
\[
\lambda_1=\lambda_2=-\lambda_7=-\mu_2=\mu_7=\mu_8=1,
\]
and all remaining coefficients equal to zero. Thus, \eqref{eq_lemma1} and \eqref{eq_lemma2} are satisfied.
\end{example}

We observe that the results of \Cref{lemma_unit} do not hold for a variety $\V$ that is not unit-closed. For instance, if $\V = \Lie$, then $\E(X) \cong \Der(X)$. However, the identity map $\id_X$ is a derivation of $X$ if and only if $X$ is abelian. 

Moreover, the conditions \eqref{eq_lemma1}--\eqref{eq_lemma2} are not sufficient to ensure that a variety of non-associative algebras $\V$ satisfying no identities of degree~$2$ is unit-closed. For instance, let $\V = \Nil_3(\Assoc)$ and fix the choice of the $\lambda/\mu$-rules as in \Cref{Example associative}. Then \eqref{eq_lemma1}--\eqref{eq_lemma2} are satisfied, but $\V$ is not unit-closed. Nevertheless, the following characterisation holds.

\begin{corollary}
Let $\V$ be an operadic, action accessible variety of non-associative algebras, and fix a choice of the $\lambda/\mu$-rules. Suppose moreover that $\V$ satisfies no identities of degree $2$ and that all the identities of $\V$ are consequences of the $\lambda/\mu$-rules. Then $\V$ is unit-closed if and only if \eqref{eq_lemma1}--\eqref{eq_lemma2} hold.
\end{corollary}

\begin{proof}
If $\V$ is unit-closed, then \eqref{eq_lemma1}--\eqref{eq_lemma2} hold by \Cref{lemma_unit}.

Conversely, assume that \eqref{eq_lemma1}--\eqref{eq_lemma2} hold. We must show that $\V$ is unit-closed, i.e., for any algebra $X$ in $\V$, the algebra $\widetilde X=\langle X,1 \rangle$ is still an object of~$\V$.

Since every identity of $\V$ is a consequence of the $\lambda/\mu$-rules, it is sufficient to verify that \eqref{eq00}--\eqref{eq01} hold for all $x,y,z \in \widetilde X$. These identities already hold when $x,y,z \in X$, since $X$ is an algebra of $\V$. Thus, it remains to check them in the cases where at least one of the variables is the external unit $1$. That is, it is enough to consider triples of the form
\[
(1,y,z),\ (x,1,z),\ (x,y,1),\ (1,1,z),\ (1,y,1),\ (x,1,1),\ (1,1,1).
\]
In each of these cases, the identities \eqref{eq00}--\eqref{eq01} reduce to trivial equalities as a consequence of \eqref{eq_lemma1}--\eqref{eq_lemma2}. This completes the proof.
\end{proof}

We are now ready to prove our main result, namely that, given an operadic, unit-closed variety $\V$, the existence of $\lambda/\mu$-rules is sufficient to guarantee that the ideally exact category $\V_1$ is action representable. 

We begin by observing that, for a homogeneous, unit-closed variety $\V$, the existence of the $\lambda/\mu$-rules in $\V_1$ is enough to ensure the existence of the $\lambda/\mu$-rules in~$\V$. Indeed, assume that there exist $\lambda_1,\ldots,\lambda_8,\mu_1,\ldots,\mu_8 \in \F$ such that the $\lambda/\mu$-rules hold in $\V_1$. Then, any (not necessarily unital) algebra of $X$ of $\V$ can be identified with the kernel of the projection
\[
\pi_1 \colon \F \ltimes X \to \F \colon (n,x) \mapsto n.
\]
Hence, $X$ is (isomorphic to) an ideal of the unital algebra $\F \ltimes X$, and therefore the $\lambda/\mu$-rules hold on $X$ as well. It follows that $\V$ satisfies the $\lambda/\mu$-rules, and we have the following.

\begin{proposition}
Let $\V$ be a homogeneous, unit-closed variety of non-associative algebras. The following are equivalent:
\begin{tfae}
\item $\V$ is action accessible;
\item The $\lambda/\mu$-rules hold in $\V$;
\item The $\lambda/\mu$-rules hold in $\V_1$. \noproof
\end{tfae}
\end{proposition}

To conclude, we have the following.

\begin{theorem}\label{thm_unital}
Let $\V$ be an operadic, action accessible, unit-closed variety of non-associative algebras, and fix a choice of the $\lambda/\mu$-rules. Let $X$ be an object of $\V$ and let $\widetilde{U} \colon \V \to \PAlg$ be the forgetful functor. The map
\[
\Inn \colon \widetilde{U}(X) \to \E(X) \colon x \mapsto (L_x,R_x)
\]
is an algebra isomorphism if and only if $X$ is a unital algebra.
\end{theorem}

\begin{proof}
If $\Inn$ is an isomorphism, then $X$ is a unital algebra with unit
\[
1=\Inn^{-1}(\id_X,\id_X).
\]
Conversely, suppose that $X$ has a multiplicative unit $1$. If $\V$ satisfies an identity of degree $2$, then by (1) of \Cref{lemma_unit}, either $\V=\CAssoc$, or $\bchar(\F)=3$ and $\V$ is a subvariety of $\JJord$. Thus, by Examples~\ref{ex_assoc} and~\ref{ex_jjord}, $\Inn$ is an isomorphism. 

Now, suppose that $\V$ does not satisfy any identity of degree $2$. Since $X$ is unital, then $\Ker(\Inn)=\Ann(X)=\{ 0 \}$, and $\Inn$ is injective. It remains to show that $\Inn$ is surjective. To this end, given any $f=(f \ast -, - \ast f) \in \E(X)$ and $x \in X$, we apply the identity \eqref{eq00} of the $\lambda/\mu$-rules with $(x,y,z)=(x,f,1)$, where as before $fx=f \ast x$ and $xf=x \ast f$. We thus have
\[
x(f \ast 1)=(\lambda_1+\lambda_3+\lambda_5+\lambda_6)(x \ast f)+(\lambda_2+\lambda_4+\lambda_7+\lambda_8)(f \ast x)
\]
and, by (2) of \Cref{lemma_unit}, it reduces to
\[
x \ast f=x(f \ast 1).
\]
Moreover, for $x=1$, we obtain $f \ast 1=1 \ast f$. Similarly, the identity \eqref{eq01} with $(x,y,z)=(x,1,f)$ gives
\[
f \ast x=(f \ast 1)x.
\]
This implies that
\[
f=(L_a,R_a)=\Inn(a)
\]
where $a=f \ast 1 = 1 \ast f \in X$, and $\Inn$ is surjective. We conclude that the map $\Inn$ is an isomorphism.
\end{proof}

As a consequence, we get the following.

\begin{theorem}\label{final_thm}
Let $\V$ be an operadic, action accessible, unit-closed variety of non-associative algebras. The ideally exact category $\V_1$ of unital algebras of $\V$ is action representable, and the actor of a unital algebra $X$ is $X$ itself. \noproof
\end{theorem}

\begin{proof}
It follows from \Cref{prop_e(x)} and \Cref{thm_unital} that, for any unital algebra $X$ of $\V$, there exists a natural isomorphism
\[
\SplExt(-,U(X)) \cong \Hom_\V(-,U(X)),
\]
where $U \colon \V_1 \rightarrow \V$ is the forgetful functor. Thus, $\V_1$ is action representable, and $X$ is its own actor.
\end{proof}

\begin{examples}{\ }
\begin{enumerate}
\item The ideally exact categories $\Assoc_1$ and $\CAssoc_1$ of unital associative and unital commutative associative algebras are action representable. Indeed, as shown in \Cref{ex_assoc}, for any unital associative algebra $X$ (resp.~commutative associative algebra), the map $\Inn$ induces an isomorphism $\Bim(X) \cong X$ (resp.~$\M(X) \cong X$).

\item
Let $\bchar(\F) \neq 2$ and consider the variety $\Alt$ of alternative algebras. Then, the ideally exact category $\Alt_1$ of unital alternative algebras is action representable. In fact, as observed in~\cite[Example 3.11(8)]{WRAAlg}, for any unital alternative algebra $X$, there is a natural isomorphism
\[
\SplExt(-,U(X)) \cong \Hom_\Alt(-,U(X)).
\]

\item Let $\WAlt$ be the variety determined by the identity
\begin{equation}\label{xabi}
x(yz)-(xy)z - z(xy) + (zx)y = 0.
\end{equation}
We call $\WAlt$ the variety of \emph{weak alternative algebras}. One may check that a choice of the $\lambda/\mu$-rules is given by
\[
\lambda_1=\lambda_3=-\lambda_6=-\mu_3=\mu_6=\mu_8=1,
\]
and hence $\WAlt$ is an operadic, action accessible, unit-closed variety.

Now, the identity \eqref{xabi} can be rewritten as
\[
[x,y,z]-[z,x,y]=0,
\]
where $[x,y,z]=x(yz)-(xy)z$ denotes the \emph{associator}. It is well known that the associator of any alternative algebra $X$ is totally skew-symmetric, that is
\[
[x_{\sigma(1)},x_{\sigma(2)},x_{\sigma(3)}]=\operatorname{sgn}(\sigma)[x_1,x_2,x_3],
\]
for any permutation $\sigma \in S_3$. In particular, since the cycle $(1\,3\,2)$ is even, we obtain
\[
[z,x,y]=[x,y,z].
\]
Therefore, every alternative algebra satisfies \eqref{xabi}, and hence $\Alt \subseteq \WAlt$. Note that this inclusion still holds in characteristic $2$.

Now, by \Cref{final_thm}, the ideally exact category $\WAlt_1$ is action representable, and the actor of any algebra $X$ of $\WAlt_1$ is $X$ itself. Consequently, the category of unital alternative algebras in characteristic $2$ is action representable too.

\item Let $\bchar(\F)=3$ and consider the variety $\JJord$. Then, the ideally exact category $\JJord_1$ of unital Jacobi-Jordan algebras is action representable. In fact, as observed in \Cref{ex_jjord}, the partial algebra $\ADer(X)$ of anti-derivations of $X$ is isomorphic to $X$ precisely when $X$ is unital.
\end{enumerate}
\end{examples}

The result of \Cref{final_thm} highlights a contrast with the non-unital case. Indeed, in the setting of semi-abelian varieties of non-associative algebras, as shown in~\cite{Tim}, action representability is a highly restrictive property, essentially characterising the variety of Lie algebras. Here, actions on $X$ become representable whenever $X$ is a unital algebra of an operadic, action accessible, unit-closed variety.

\section{Action representability of unital Poisson algebras}\label{SecPois}

The aim of this section is to study the representability of actions of unital Poisson algebras. We prove that the categories $\Pois_1$ and $\CPois_1$ of unital Poisson algebras and unital commutative Poisson algebras over a field $\F$ with $\bchar(\F) \neq 2$ are action representable, by employing the explicit construction of the universal strict general actor given in~\cite{CigoliManciniMetere}. This also suggests a possible approach to extending the present analysis to ideally exact categories of unital algebras with two bilinear, not necessarily associative, operations.

\begin{definition}
A Poisson algebra is a vector space $X$ over a field $\F$ equipped with two bilinear operations
\[
\cdot \colon X \times X\to X \quad \text{and} \quad [-,-] \colon X \times X  \to X,
\]
such that $(X,\cdot)$ is an associative algebra, $(X,[-,-])$ is a Lie algebra and the \emph{Poisson identity} holds:
\[
[x,yz]=[x,y]z+y[x,z], \quad \forall x,y,z \in X.
\]
This means that for any $x \in X$, the adjoint map $\operatorname{ad}_x=[x,-]$ is a derivation of the associative algebra $(X,\cdot)$.

A Poisson algebra is said to be \emph{commutative} (resp.~\emph{unital}) if the underlying associative algebra is commutative (resp.~unital).
\end{definition}

\begin{remark}
Let $X$ be a Poisson algebra. We denote by
\[
\Ann(X)=\{ x \in X \mid xy=yx=0, \; \forall y \in X \}
\]
the annihilator of the associative algebra $(X,\cdot)$, and by
\[
\bZ(X)=\{ x \in X \mid [x,y]=0, \; \forall y \in X \}
\]
the center of the Lie algebra $(X,[-,-])$.
\end{remark}

We observe that, if $X$ is a unital Poisson algebra, then the unit element $1$ belongs to the center $\bZ(X)$. Indeed, applying the Poisson identity with $y=z=1$, one gets
\[
[x,1]=[x,1]1+1[x,1]=2[x,1]
\]
and, hence, $[x,1]=[1,x]=0$ for any $x \in X$.

\begin{example}
Any associative algebra $(X,\cdot)$ becomes a Poisson algebra with the usual Lie bracket
\[
[x,y]=xy-yx.
\]
Furthermore:
\begin{itemize}
\item If $(X,\cdot)$ is a commutative associative algebra, then $(X,\cdot,[-,-])$ is a commutative Poisson algebra with trivial Lie bracket.
\item If $(X,\cdot)$ is a unital associative algebra, then $(X,\cdot,[-,-])$ is a unital Poisson algebra.
\end{itemize}
\end{example}

We denote by $\Pois$ and $\CPois$ the semi-abelian varieties of (commutative) Poisson algebras, and by $\Pois_1$ and $\CPois_1$ the subcategories of unital (commutative) Poisson algebras. We observe that $\Pois$ and $\CPois$ are Orzech categories of interest.

\begin{remark}
The categories $\Pois_1$ and $\CPois_1$ are ideally exact, and their initial object is the field $\F$ endowed with the standard multiplication and the trivial Lie bracket $[n,m]=0$. Furthermore, there is a monadic adjunction
\[
\begin{tikzcd}
{\Pois_1} & {\Pois,}
\arrow[""{name=0, anchor=center, inner sep=0}, "U"', from=1-1, to=1-2]
\arrow[""{name=1, anchor=center, inner sep=0}, "F"', curve={height=12pt}, from=1-2, to=1-1]
\arrow["\dashv"{anchor=center, rotate=-90}, draw=none, from=1, to=0]
\end{tikzcd}
\]
where $U$ is the forgetful functor and $F$ maps every not necessarily unital Poisson algebra $X$ to the semidirect product $\F \ltimes X$ with operations
\begin{align*}
(n,x)+(m,y)&=(n+m,x+y),\\
(n,x)\cdot(m,y)&=(nm, xy + ny + mx),\\
[(n,x),(m,y)]&=(0,[x,y])
\end{align*}
and unit $(1,0)$. As in the case of associative algebras, the unit $\eta \colon 1_{\Pois} \Rightarrow UF$ of the adjunction is cartesian since the map $\eta_X \colon X \to  U(\F \ltimes X) \colon x \mapsto (0,x)$ is the kernel of
\[
U(\pi_1) \colon U(\F \ltimes X) \to  U(\F) \colon (n,x) \mapsto n.
\]
Thus, $F \dashv U$ is, up to an equivalence, the adjunction associated with the unique morphism $\F \to \{ 0 \}$ in $\Pois_1$.

Similarly, there is a monadic adjunction with cartesian unit
\[
\begin{tikzcd}
{\CPois_1} & {\CPois.}
\arrow[""{name=0, anchor=center, inner sep=0}, "U"', from=1-1, to=1-2]
\arrow[""{name=1, anchor=center, inner sep=0}, "F"', curve={height=12pt}, from=1-2, to=1-1]
\arrow["\dashv"{anchor=center, rotate=-90}, draw=none, from=1, to=0]
\end{tikzcd}
\]
\end{remark}

Since $\Pois$ is an Orzech category of interest, it follows from~\cite[Corollary 4.2]{CigoliManciniMetere} that, for every Poisson algebra $X$, there exists a natural monomorphism of functors
\[
\tau \colon \SplExt(-,X) \rightarrowtail \Hom_{\Alg^2}(\overline{U}(-),\USGA(X)).
\]
Here, $\Alg^2$ denotes the category of algebras with two not necessarily associative bilinear operations, $\overline{U} \colon \Pois \to \Alg^2$ is the forgetful functor, and $\USGA(X)$ is the universal strict general actor of $X$ (see~\cite{Casas}).

It is shown in~\cite[Theorem 5.6]{CigoliManciniMetere} that $\USGA(X)$ can be described as the vector space
\begin{align*}
\USGA(X)=\{ f=(f \ast -,- \ast f,[f,-]) &\in \Bim(X)\times \Der(X) \mid \cdots \\
&\quad \cdots \mid f \ast [x,y]=[f \ast x,y]-[f,y]x, \\
&\qquad\quad [x,y] \ast f=[x \ast f,y]-x[f,y], \\
&\qquad\quad [f,xy]=[f,x]y+x[f,y], \; \forall x,y \in X \},
\end{align*}
endowed with the bilinear operations
\[
f \cdot g=(f \ast (g \ast -), (- \ast f) \ast g), f \ast [g,-] + [f,-] \ast g)
\]
and
\[
[f,g]=(f \ast [g,-] - [g,f \ast -], [g,-] \ast f - [g,- \ast f], [f,[g,-]]-[g,[f,-]]).
\]

Moreover, a morphism 
\[
\varphi \colon \overline{U}(B) \to \USGA(X) \colon b \mapsto (b \ast -, - \ast b, [b,-])
\]
in $\Alg_2$ belongs to $\Imm(\tau_B)$ if and only if 
\begin{equation}\label{eq_pois}
(b \ast x) \ast b' = b \ast (x \ast b'),
\end{equation}
for any $b,b' \in B$ and $x \in X$ (see~\cite[Theorem 5.6]{CigoliManciniMetere}).

\begin{remark}
Let $X$ be a Poisson algebra. We observe that the non-associative algebra $(\USGA(X),\cdot)$ has multiplicative unit
\[
1=(\id_X,\id_X,0_X),
\]
where $0_X \colon X \to X$ denotes the zero map. Indeed, one has
\[
f \cdot 1=1 \cdot f=f,
\]
for any $f \in \USGA(X)$. Furthermore, $1 \in \bZ(\USGA(X))$, since $[f,1]=[1,f]=0_X$.
\end{remark}

We further recall from~\cite[Remark 5.8]{CigoliManciniMetere} that, for any Poisson algebra $X$, if the underlying associative algebra $(X,\cdot)$ has trivial annihilator or is perfect, that is, $(X^2,\cdot)=(X,\cdot)$, then
\[
f \ast (x \ast g) = (f \ast x) \ast g,
\]
for every $f=(f \ast -, - \ast f)$ and $g=(g \ast -,- \ast g)$ in $\Bim(X)$ and $x \in X$.

Consequently, for any other Poisson algebra $B$, every morphism $\overline{U}(B) \to \USGA(X)$ lies in $\Imm(\tau_B)$, and hence there is a natural isomorphism
\begin{equation}\label{isopois}
\SplExt(-,X)\cong \Hom_{\Alg^{2}}(\overline{U}(-),\USGA(X)).
\end{equation}

Since any unital associative algebra is perfect and with trivial annihilator, the isomorphism \eqref{isopois} holds in particular for any object $X$ of $\Pois_1$. Moreover, we can prove the following.

\begin{proposition}\label{prop_pois}
Let $X$ be a Poisson algebra. The algebra homomorphism
\[
\Inn \colon \overline{U}(X) \to \USGA(X) \colon x \mapsto (L_x, R_x, \operatorname{ad}_x),
\] 
is an isomorphism if and only if $X$ is a unital algebra.
\end{proposition}

\begin{proof}
If $\Inn$ is an isomorphism, then $X$ is a unital Poisson algebra with
\[
1=\Inn^{-1}(\id_X,\id_X,0_X).
\]
Conversely, suppose that $X$ is a unital Poisson algebra. We have
\[
\Ker(\Inn)=\Ann(X)=\{0\},
\]
and $\Inn$ is injective. It remains to show that $\Inn$ is surjective. To this end, let
\[
f=(f \ast -, - \ast f, [f,-]) \in \USGA(X).
\]
Since $(X,\cdot)$ is unital, by \Cref{ex_assoc} we have an isomorphism of associative algebras $\Bim(X) \cong X$, and there exists $a \in X$ such that
\[
f \ast x= ax \quad \text{and} \quad x \ast f= xa
\]
for any $x \in X$. Moreover, since $a=f \ast 1=1 \ast f$, one has
\[
0=f \ast [1,x]=[f \ast 1,x]-[f,x]1=[a,x]-[f,x],
\]
and hence $[f,x]=[a,x]$ for every $x \in X$. Thus, $f=(L_a,R_a,\operatorname{ad}_a)=\Inn(a)$, and we conclude that $\Inn$ is an isomorphism.
\end{proof}

It follows that for any unital Poisson algebra $X$, there is a natural isomorphism
\[
\SplExt(-, U(X)) \cong \Hom_\Pois(-,U(X)),
\]
where $U \colon \Pois_{1} \to \Pois$ denotes the forgetful functor, and we get the following.

\begin{theorem}
The ideally exact category $\Pois_{1}$ is action representable, and the actor of a unital Poisson algebra $X$ is $X$ itself. \noproof
\end{theorem}

We observe that the same result may be reached in the category $\CPois_{1}$ of unital commutative Poisson algebras, where an easier description of the universal strict general actor as a subspace of $\M(X) \times \Der(X)$ is available (see~\cite{CigoliManciniMetere}). 
Hence, we conclude with the following.

\begin{theorem}
The category $\CPois_{1}$ is action representable, and the actor of a unital commutative Poisson algebra $X$ is $X$ itself. \noproof
\end{theorem}

\section*{Acknowledgements}
This work is supported by the University of Messina, University of Palermo and by the ``National Group for Algebraic and Geometric Structures and their Applications'' (GNSAGA -- INdAM). The first author is supported by the SDF Sustainability Decision Framework Research Project -- MISE decree of 31/12/2021 (MIMIT Dipartimento per le politiche per le imprese -- Direzione generale per gli incentivi alle imprese) -- CUP:~B79J23000530005, COR:~14019279, Lead Partner:~TD Group Italia Srl, Partner:~University of Palermo, and he is also a Postdoctoral Researcher of the Fonds de la Recherche Scientifique--FNRS. The third author is supported by Wallonie-Bruxelles International.

%\section*{Declarations}

%\subsection*{Corresponding Author}
%The first author is the corresponding author.

%\subsection*{Availability of data and materials}
%Not applicable.

%\subsection*{Competing interests}
%Not applicable. There is no competing interest.

%\subsection*{Authors' contributions}
%All authors contributed equally to this article.

%\section*{Acknowledgements}
%There is no Competing Interest.

\begin{thebibliography}{99}

\bibitem{Barr}
M.~Barr, \emph{Exact Categories}, Lecture Notes in Mathematics \textbf{236} (1971), 1--120.

\bibitem{Bell}
J.~Bell and M.~Machover, \emph{A course in mathematical logic}, North-Holland, 1977.

\bibitem{malcev}
F.~Borceux and D.~Bourn, \emph{Mal'cev, Protomodular, Homological and Semi-Abelian Categories}, Kluwer Academic Publishers, 2004. \doi{10.1007/978-1-4020-1962-3}

\bibitem{BJK}
F.~Borceux, G.~Janelidze and G.~M. Kelly, \emph{Internal object actions}, Commentationes Mathematicae Universitatis Carolinae \textbf{46} (2005), no.~2, 235--255.

\bibitem{BJK2}
F.~Borceux, G.~Janelidze and G.~M. Kelly, \emph{On the representability of actions in a semi-abelian category}, Theory and Applications of Categories \textbf{14} (2005), no.~11, 244--286.

\bibitem{Bourn}
D.~Bourn, \emph{Normalization equivalence, kernel equivalence and affine categories}, Lecture Notes in Mathematics \textbf{1488} (1991), 43--62.

\bibitem{act_accessible}
D.~Bourn and G.~Janelidze, \emph{Centralizers in action accessible categories}, Cahiers de Topologie et G\'eom\'etrie Diff\'erentielle Cat\'egoriques \textbf{50} (2009), no.~3, 211--232.

\bibitem{WRAAlg}
J.~Brox, X.~García-Martínez, M.~Mancini, T.~Van~der Linden and C.~Vienne, \emph{Weak representability of actions of non-associative algebras}, Journal of Algebra \textbf{669} (2025), no.~18, 401--444. \doi{10.1016/j.jalgebra.2025.02.007}

\bibitem{Burde}
D.~Burde, A.~Fialowski, \emph{Jacobi–Jordan algebras}, Linear Algebra and its Applications \textbf{459} (2014), no. 34, 586--594. \doi{10.1016/j.laa.2014.07.034}

\bibitem{Cappelletti}
A.~Cappelletti, \emph{Categorical-Algebraic Properties of Lattice-ordered Groups}, Theory and Applications of Categories \textbf{39} (2023), no.~31, 916--948.

\bibitem{Casas}
J.~M.~Casas, T.~Datuashvili and M.~Ladra, \emph{Universal strict general actors and actors in categories of interest}, Applied Categorical Structures \textbf{18} (2010), 85--114. \doi{10.1007/s10485-008-9166-z}

\bibitem{CigoliManciniMetere}
A.~S. Cigoli, M.~Mancini and G.~Metere, \emph{On the representability of actions of Leibniz algebras and Poisson algebras}, Proceedings of the Edinburgh Mathematical Society \textbf{66} (2023), no.~4, 998--1021. \doi{10.1017/S0013091523000548}

\bibitem{XabiMancini}
X.~García-Martínez and M.~Mancini, \emph{Action accessible and weakly action representable varieties of algebras}, Bollettino dell'Unione Matematica Italiana (2025). \doi{10.1007/s40574-025-00496-1}

\bibitem{Tim}
X.~García-Martínez, M.~Tsishyn, T.~Van der Linden and C.~Vienne, \emph{Algebras with representable representations}, Proceedings of the Edinburgh Mathematical Society \textbf{64} (2021), no.~2, 555--573. \doi{10.1017/S0013091521000304}

\bibitem{Gray}
J.~R.~A. Gray, \emph{A note on the relationship between action accessible and weakly action representable categories}, Theory and Applications of Categories \textbf{44} (2025), no.~8, 272--276.

\bibitem{WRA}
G.~Janelidze, \emph{Central extensions of associative algebras and weakly action representable categories}, Theory and Applications of Categories \textbf{38} (2022), no.~36, 1395--1408.

\bibitem{IdeallyExact}
G.~Janelidze, \emph{Ideally exact categories}, Theory and Applications of Categories \textbf{41} (2024), no.~11, 414--425.

\bibitem{Semi-Ab}
G.~Janelidze, L.~Márki and W.~Tholen, \emph{Semi-abelian categories}, Journal of Pure and Applied Algebra \textbf{168} (2002), no.~2, 367--386. \doi{10.1016/S0022-4049(01)00103-7}

\bibitem{rel}
S.~Lapenta, G.~Metere and L.~Spada, \emph{Relative ideals in homological categories, with an application to MV-algebras}, Theory and Applications of Categories \textbf{41} (2024), no.~27, 878--893.

\bibitem{loday}
J.-L.~Loday, \emph{Une version non commutative des algèbres de Lie: les algèbres de Leibniz}, L'Enseignement Math\'ematique \textbf{39} (1993), no.~3-4, 269--293.

\bibitem{MacLane58}
S.~Mac~Lane, \emph{Extensions and obstructions for rings}, Illinois Journal of Mathematics \textbf{2} (1958), no.~3, 316--345. \doi{10.1215/ijm/1255454537}

\bibitem{Mancini_Thesis}
M.~Mancini, \emph{On the representability of actions of non-associative algebras}, Ph.D. Thesis (2024), Università degli Studi di Palermo, Italy. \url{https://hdl.handle.net/10447/623093}

\bibitem{Montoli}
A.~Montoli, \emph{Action accessibility for categories of interest}, Theory and Applications of Categories \textbf{23} (2010), no.~1, 7--21.

\bibitem{Orzech} 
G.~Orzech, \emph{Obstruction theory in algebraic categories {I} and {II}}, Journal of Pure and Applied Algebra~\textbf{2} (1972), no.~4, 287--314 and 315--340. \doi{10.1016/0022-4049(72)90008-4}

\bibitem{Osborn}
J.~M.~Osborn, \emph{Varieties of algebras}, Advances in Mathematics~\textbf{8} (1972), 163--369.\\ \doi{10.1016/0001-8708(72)90003-5}

\bibitem{VdL-NAA}
T.~Van~der Linden, \emph{Non-associative algebras}, New Perspectives in Algebra, Topology and Categories (M.~M. Clementino, A.~Facchini, M.~Gran, eds.), Coimbra Mathematical Texts, vol.~1, Springer, Cham, 2021, 225--258. \doi{\detokenize{10.1007/978-3-030-84319-9_7}}

\bibitem{Coco_Thesis}
C.~Vienne, \emph{Categorical-algebraic conditions in semi-abelian categories}, Ph.D. Thesis (2025), Université catholique de Louvain, Belgium. \url{http://hdl.handle.net/2078.1/301231}

\bibitem{hmg}
K.~A.~Zhevlakov, A.~M.~Slin’ko, I.~P.~Shestakov and A.~I.~Shirshov, \emph{Rings that are nearly associative} (1982), Academic Press. 

\end{thebibliography}
\end{document}
